\documentclass{article}
\usepackage[utf8]{inputenc}
\usepackage{amsmath, amssymb, graphicx, xcolor, fancyvrb}

% --- UBAH JUDUL, NAMA, DAN NIM ANDA DI SINI ---
\title{TugasAplikom4Adi25030630001}
\author{Adi Kurniawan \\ NIM: 25030630001}
\date{\today}

% --- PENERJEMAH OTOMATIS KODE EMT (JANGAN DIUBAH) ---
\newenvironment{eulernotebook}{\par}{\par}
\newcommand{\eulerheading}[1]{\section{#1}}
\newcommand{\eulersubheading}[1]{\subsection{#1}}
\newenvironment{eulercomment}{\par\medskip}{\par\medskip}
\newenvironment{eulerformula}{}{}
\DefineVerbatimEnvironment{eulerprompt}{Verbatim}{formatcom=\color{red}}
\DefineVerbatimEnvironment{euleroutput}{Verbatim}{formatcom=\color{blue}}
\DefineVerbatimEnvironment{eulerttcomment}{Verbatim}{formatcom=\color{teal}}
\newcommand{\eulerimg}[2]{\begin{center}\includegraphics[width=0.8\textwidth]{#2}\end{center}}

\begin{document}
\maketitle

% ====================================================================
% TEMPELKAN KODE DARI \begin{eulernotebook} SAMPAI \end{eulernotebook}
% DI BAWAH BARIS INI:
\begin{eulernotebook}
\eulerheading{Menggambar Grafik 2D dengan EMT}
\begin{eulercomment}
Notebook ini menjelaskan tentang cara menggambar berbagaikurva dan
grafik 2D dengan software EMT. EMT menyediakan fungsi plot2d() untuk
menggambar berbagai kurva dan grafik dua dimensi (2D).

\end{eulercomment}
\eulersubheading{Basic Plots}
\begin{eulercomment}
Ada beberapa fungsi paling dasar dari grafik. Ada koordinat layar,
yang nilainya selalu berkisar dari 0 sampai 1024 di setiap sumbu,
nggak peduli bentuk layarnya persegi atau bukan. Lalu ada koordinat
grafik, yang bisa diatur pake setplot(). Pemetaan antara kedua
koordinat ini tergantung pada jendela grafik yang lagi aktif.
Contohnya, fungsi bawaan shrinkwindow() bakal nyediain ruang sisanya
buat label sumbu dan judul grafik.

Di contoh ini, kita cuma menggambar beberapa garis acak dengan
berbagai warna. Buat detail fungsi-fungsi ini, kamu bisa pelajari
fungsi inti dari EMT.
\end{eulercomment}
\begin{eulerprompt}
>clg; // clear screen
>window(0,0,1024,1024); // use all of the window
>setplot(0,1,0,1); // set plot coordinates
>hold on; // start overwrite mode
>n=100; X=random(n,2); Y=random(n,2);  // get random points
>colors=rgb(random(n),random(n),random(n)); // get random colors
>loop 1 to n; color(colors[#]); plot(X[#],Y[#]); end; // plot
>hold off; // end overwrite mode
>insimg; // insert to notebook
\end{eulerprompt}
\eulerimg{27}{Images Use/TugasAplikom4Adi25030630001-001.png}
\begin{eulerprompt}
>reset;
\end{eulerprompt}
\begin{eulercomment}
Grafiknya perlu ditahan dulu (di-hold), soalnya perintah plot() bakal
menghapus jendela grafik yang ada.

Buat menghapus semua yang udah kita buat, kita pake reset().

Untuk menampilkan gambar hasil plot di layar notebook, perintah
plot2d() dapat diakhiri dengan titik dua (:). Cara lain adalah
perintah plot2d() diakhiri dengan titik koma (;), kemudian menggunakan
perintah insimg() untuk menampilkan gambar hasil plot.

Sebagai contoh lain, kita menyisipkan sebuah grafik ke dalam grafik
lain (inset). Ini dilakukan dengan mendefinisikan jendela grafik yang
lebih kecil. Perhatikan bahwa jendela ini tidak menyediakan ruang
untuk label sumbu di luar area grafik. Kita perlu menambahkan sedikit
margin untuk itu sesuai kebutuhan. Perhatikan juga bahwa kita
menyimpan dan mengembalikan tampilan jendela penuh (full window),
serta menahan (hold) grafik yang sedang aktif saat kita menggambar
grafik sisipan tersebut.
\end{eulercomment}
\begin{eulerprompt}
>reset()
\end{eulerprompt}
\begin{euleroutput}
  0
\end{euleroutput}
\begin{eulerprompt}
>plot2d("x^3-x");
>plot2d("x^3-x");
>insimg()
\end{eulerprompt}
\eulerimg{27}{Images Use/TugasAplikom4Adi25030630001-002.png}
\begin{eulerprompt}
>xw=200; yw=100; ww=300; hw=300;
>ow=window();
>window(xw,yw,xw+ww,yw+hw);
>hold on;
>barclear(xw-50,yw-10,ww+60,ww+60);
>plot2d("x^4-x",grid=6):
\end{eulerprompt}
\eulerimg{27}{Images Use/TugasAplikom4Adi25030630001-003.png}
\begin{eulerprompt}
>hold off;
>window(ow);
>figure(2,2); ...
>for n=1 to 4; figure(n); plot2d("x^"+n); end; ...
>figure(0):
\end{eulerprompt}
\eulerimg{27}{Images Use/TugasAplikom4Adi25030630001-004.png}
\begin{eulercomment}
Untuk membuat grafik dengan beberapa gambar sekaligus, caranya sama
saja. Ada fungsi bantuan figure() yang bisa dipakai untuk hal ini.
\end{eulercomment}
\begin{eulercomment}

\end{eulercomment}
\eulersubheading{Aspek Plot}
\begin{eulercomment}
Tampilan grafik bawaan menggunakan jendela grafik berbentuk persegi.
Kamu bisa mengubahnya dengan fungsi aspect(). Jangan lupa untuk
mengembalikan (reset) pengaturan aspeknya nanti. Kamu juga bisa
mengubah nilai bawaan ini lewat menu pada opsi "Set Aspect" ke rasio
aspek tertentu atau disesuaikan dengan ukuran jendela grafik saat ini.

Tapi kamu juga bisa mengubahnya khusus untuk satu grafik saja.
Caranya, ukuran area grafik yang sedang aktif akan diubah, lalu
jendelanya disesuaikan supaya label grafik punya cukup ruang.
\end{eulercomment}
\begin{eulerprompt}
>aspect(2); // rasio panjang dan lebar 2:1
>plot2d(["sin(x)","cos(x)"],0,2pi):
\end{eulerprompt}
\eulerimg{13}{Images Use/TugasAplikom4Adi25030630001-005.png}
\begin{eulerprompt}
>aspect();
>reset;
\end{eulerprompt}
\begin{eulercomment}
Fungsi reset() mengembalikan pengaturan bawaan grafik, termasuk rasio
aspek (aspect ratio).\\
\begin{eulercomment}
\eulerheading{2D Plots in Euler}
\begin{eulercomment}
EMT Math Toolbox memiliki fitur grafik 2D, baik untuk data maupun
fungsi. EMT menggunakan fungsi plot2d. Fungsi ini dapat memplot fungsi
dan data.

Kalian juga bisa membuat grafik di Maxima menggunakan Gnuplot atau di
Python menggunakan Matplotlib.

Euler dapat membuat grafik 2D untuk:

- ekspresi,\\
fungsi, variabel, atau kurva terparameterisasi,

-vektor nilai x-y,

-kumpulan titik (point cloud) pada bidang datar,

-kurva implisit dengan level atau wilayah level (level regions),

-fungsi kompleks.

Gaya grafik mencakup berbagai opsi gaya untuk garis dan titik, grafik
batang, serta grafik berarsir (shaded plots).\\
\begin{eulercomment}
\eulerheading{Plot Ekspresi atau Variabel}
\begin{eulercomment}
Sebuah ekspresi tunggal dalam variabal "x" (misalnya "4*x\textasciicircum{}2") atau
nama dari suatu fungsi (misalnya "f") akan menghasilkan grafik dari
fungsi tersebut.

Berikut adalah contoh paling dasar, yang menggunakan rentang bawaan
(default range) dan menyesuaikan rentang sumbu-y agar cocok dengan
grafik fungsi tersebut.

Catatan: Jika kamu mengakhiri baris perintah dengan tanda titik dua
":", grafik akan langsung disisipkan ke dalam jendela teks. Jika
tidak, tekan tombol TAB untuk melihat grafik apabila jendela grafik
tertutup.
\end{eulercomment}
\begin{eulerprompt}
>plot2d("x^2"):
\end{eulerprompt}
\eulerimg{27}{Images Use/TugasAplikom4Adi25030630001-006.png}
\begin{eulerprompt}
>aspect(1.5); plot2d("x^3-x"):
\end{eulerprompt}
\eulerimg{17}{Images Use/TugasAplikom4Adi25030630001-007.png}
\begin{eulerprompt}
>a:=5.6; plot2d("exp(-a*x^2)/a"); insimg(30); // menampilkan gambar hasil plot setinggi 25 baris
\end{eulerprompt}
\eulerimg{17}{Images Use/TugasAplikom4Adi25030630001-008.png}
\begin{eulercomment}
Dari beberapa contoh sebelumnya Anda dapat melihat bahwa aslinya
gambar plot menggunakan sumbu X dengan rentang nilai dari -2 sampai
dengan 2. Untuk mengubah rentang nilai X dan Y, Anda dapat menambahkan
nilai-nilai batas X (dan Y) di belakang ekspresi yang digambar.

Rentang grafik (plot range) diatur menggunakan parameter-parameter
berikut:

- a,b: rentang sumbu-x (bawaan: -2, 2)\\
- c,d: rentang sumbu-y (bawaan: menyesuaikan dengan nilai data)\\
- r: opsi alternatif berupa jari-jari (radius) di sekitar titik pusat
grafik\\
- cx,cy: koordinat titik pusat grafik (bawaan: 0,0)
\end{eulercomment}
\begin{eulerprompt}
>plot2d("x^3-x",-1,2):
\end{eulerprompt}
\eulerimg{17}{Images Use/TugasAplikom4Adi25030630001-009.png}
\begin{eulerprompt}
>plot2d("sin(x)",-2*pi,2*pi): // plot sin(x) pada interval [-2pi, 2pi]
\end{eulerprompt}
\eulerimg{17}{Images Use/TugasAplikom4Adi25030630001-010.png}
\begin{eulerprompt}
>plot2d("cos(x)","sin(3*x)",xmin=0,xmax=2pi):
\end{eulerprompt}
\eulerimg{17}{Images Use/TugasAplikom4Adi25030630001-011.png}
\begin{eulercomment}
Alternatif selain tanda titik dua adalah perintah insimg(lines), yang
menyisipkan grafik dengan memakan jumlah baris teks tertentu.

Pada bagian opsi, grafik dapat diatur agar tampil:

- di jendela terpisah yang ukurannya bisa diubah,\\
- di dalam jendela notebook.

Gaya lain yang lebih bervariasi bisa didapatkan dengan menggunakan
perintah grafik khusus.

Dalam kasus apa pun, tekan tombol tabulator untuk melihat grafik jika
gambarnya tersembunyi.

Untuk membagi jendela menjadi beberapa grafik, gunakan perintah
figure(). Pada contoh ini, kita membuat grafik dari x\textasciicircum{}1 sampai x\textasciicircum{}4 ke
dalam 4 bagian jendela. figure(0) mengembalikan tampilan ke jendela
bawaan.
\end{eulercomment}
\begin{eulerprompt}
>reset;
>figure(2,2); ...
>for n=1 to 4; figure(n); plot2d("x^"+n); end; ...
>figure(0):
\end{eulerprompt}
\eulerimg{27}{Images Use/TugasAplikom4Adi25030630001-012.png}
\begin{eulercomment}
Pada plot2d(), terdapat pilihan gaya alternatif yang bisa digunakan
dengan grid=x. Untuk gambaran umumnya, kami menampilkan berbagai gaya
kisi (grid) dalam satu gambar (lihat di bawah untuk penggunaan
perintah figure()). Gaya grid=0 tidak disertakan, yang mana gaya ini
tidak menampilkan kisi maupun bingkai.
\end{eulercomment}
\begin{eulerprompt}
>figure(3,3); ...
>for k=1:9; figure(k); plot2d("x^3-x",-2,1,grid=k); end; ...
>figure(0):
\end{eulerprompt}
\eulerimg{27}{Images Use/TugasAplikom4Adi25030630001-013.png}
\begin{eulercomment}
Jika argumen untuk plot2d() berupa ekspresi yang diikuti oleh empat
angka, angka-angka tersebut merupakan rentang sumbu-x dan sumbu-y
untuk grafik tersebut.

Sebagai alternatif, parameter a, b, c, d dapat ditentukan menggunakan
nama parameter seperti a=... dan seterusnya.

Pada contoh berikut, kita mengubah gaya kisi (grid), menambahkan
label, serta menggunakan label vertikal untuk sumbu-y.
\end{eulercomment}
\begin{eulerprompt}
>aspect(1.5); plot2d("sin(x)",0,2pi,-1.2,1.2,grid=3,xl="x",yl="sin(x)"):
\end{eulerprompt}
\eulerimg{17}{Images Use/TugasAplikom4Adi25030630001-014.png}
\begin{eulerprompt}
>plot2d("sin(x)+cos(2*x)",0,4pi):
\end{eulerprompt}
\eulerimg{17}{Images Use/TugasAplikom4Adi25030630001-015.png}
\begin{eulercomment}
Gambar yang dihasilkan dengan menyisipkan grafik ke dalam jendela teks
disimpan di direktori yang sama dengan notebook, secara bawaan berada
di dalam subdirektori bernama "images". Gambar-gambar ini juga
digunakan untuk ekspor ke format HTML.

Kamu bisa menyalin gambar apa saja ke clipboard secara langsung dengan
menekan tombol Ctrl-C pada gambar yang dipilih. Tentu saja, kamu juga
bisa mengekspor grafik saat ini menggunakan fungsi-fungsi yang ada di
menu File.

Fungsi atau ekspresi pada plot2d dievaluasi secara adaptif. Untuk
kecepatan lebih tinggi, matikan grafik adaptif dengan \textless{}adaptive dan
tentukan jumlah sub-interval menggunakan n=.... Hal ini biasanya hanya
diperlukan dalam kasus-kasus tertentu yang cukup jarang.
\end{eulercomment}
\begin{eulerprompt}
>plot2d("sign(x)*exp(-x^2)",-1,1,<adaptive,n=10000):
\end{eulerprompt}
\eulerimg{17}{Images Use/TugasAplikom4Adi25030630001-016.png}
\begin{eulerprompt}
>plot2d("x^x",r=1.2,cx=1,cy=1):
\end{eulerprompt}
\eulerimg{17}{Images Use/TugasAplikom4Adi25030630001-017.png}
\begin{eulercomment}
Perhatikan bahwa \textdollar{}x\textasciicircum{}x\textdollar{} tidak terdefinisi untuk \textdollar{}x \textbackslash{}le 0\textdollar{}. Fungsi
plot2d menangani kesalahan (error) ini dan mulai menggambar grafik
segera setelah fungsinya terdefinisi. Ini berlaku untuk semua fungsi
yang mengembalikan nilai NAN di luar rentang definisinya.
\end{eulercomment}
\begin{eulerprompt}
>plot2d("log(x)",-0.1,2):
\end{eulerprompt}
\eulerimg{17}{Images Use/TugasAplikom4Adi25030630001-018.png}
\begin{eulercomment}
Parameter square=true (atau \textgreater{}square) memilih rentang sumbu-y secara
otomatis sehingga hasilnya berupa jendela grafik yang persegi.
Perhatikan bahwa secara bawaan (default), Euler menggunakan ruang
berbentuk persegi di dalam jendela grafik.
\end{eulercomment}
\begin{eulerprompt}
>plot2d("x^3-x",>square):
\end{eulerprompt}
\eulerimg{17}{Images Use/TugasAplikom4Adi25030630001-019.png}
\begin{eulerprompt}
>plot2d(''integrate("sin(x)*exp(-x^2)",0,x)'',0,2): // plot integral
\end{eulerprompt}
\eulerimg{17}{Images Use/TugasAplikom4Adi25030630001-020.png}
\begin{eulercomment}
Jika kamu membutuhkan lebih banyak ruang untuk label sumbu-y, panggil
fungsi shrinkwindow() dengan parameter smaller, atau atur nilai
positif untuk opsi smaller pada plot2d().
\end{eulercomment}
\begin{eulerprompt}
>plot2d("gamma(x)",1,10,yl="y-values",smaller=6,<vertical):
\end{eulerprompt}
\eulerimg{17}{Images Use/TugasAplikom4Adi25030630001-021.png}
\begin{eulercomment}
Ekspresi simbolik juga bisa digunakan, karena ekspresi tersebut
disimpan sebagai ekspresi string sederhana.
\end{eulercomment}
\begin{eulerprompt}
>x=linspace(0,2pi,1000); plot2d(sin(5x),cos(7x)):
\end{eulerprompt}
\eulerimg{17}{Images Use/TugasAplikom4Adi25030630001-022.png}
\begin{eulerprompt}
>a:=5.6; expr &= exp(-a*x^2)/a; // define expression
>plot2d(expr,-2,2): // plot from -2 to 2
\end{eulerprompt}
\eulerimg{17}{Images Use/TugasAplikom4Adi25030630001-023.png}
\begin{eulerprompt}
>plot2d(expr,r=1,thickness=2): // plot in a square around (0,0)
\end{eulerprompt}
\eulerimg{17}{Images Use/TugasAplikom4Adi25030630001-024.png}
\begin{eulerprompt}
>plot2d(&diff(expr,x),>add,style="--",color=red): // add another plot
\end{eulerprompt}
\eulerimg{17}{Images Use/TugasAplikom4Adi25030630001-025.png}
\begin{eulerprompt}
>plot2d(&diff(expr,x,2),a=-2,b=2,c=-2,d=1): // plot in rectangle
\end{eulerprompt}
\eulerimg{17}{Images Use/TugasAplikom4Adi25030630001-026.png}
\begin{eulerprompt}
>plot2d(&diff(expr,x),a=-2,b=2,>square): // keep plot square
\end{eulerprompt}
\eulerimg{17}{Images Use/TugasAplikom4Adi25030630001-027.png}
\begin{eulerprompt}
>plot2d("x^2",0,1,steps=1,color=red,n=10):
\end{eulerprompt}
\eulerimg{17}{Images Use/TugasAplikom4Adi25030630001-028.png}
\begin{eulerprompt}
>plot2d("x^2",>add,steps=2,color=blue,n=10):
\end{eulerprompt}
\eulerimg{17}{Images Use/TugasAplikom4Adi25030630001-029.png}
\eulerheading{Functions in one Parameter}
\begin{eulercomment}
Fungsi pemetaan (plotting) yang paling penting untuk grafik dua
dimensi (planar plots) adalah plot2d(). Fungsi ini diimplementasikan
dalam bahasa Euler di dalam berkas "plot.e", yang dimuat saat program
dijalankan.

Berikut adalah beberapa contoh penggunaan fungsi. Seperti biasa di
EMT, untuk fungsi yang bekerja pada fungsi atau ekspresi lain, kamu
bisa meneruskan parameter tambahan (selain x) yang bukan merupakan
variabel global ke fungsi tersebut menggunakan parameter titik koma
(semicolon parameters) atau dengan koleksi panggilan (call
collection).
\end{eulercomment}
\begin{eulerprompt}
>function f(x,a) := x^2/a+a*x^2-x; // define a function
>a=0.3; plot2d("f",0,1;a): // plot with a=0.3
\end{eulerprompt}
\eulerimg{17}{Images Use/TugasAplikom4Adi25030630001-030.png}
\begin{eulerprompt}
>plot2d("f",0,1;0.4): // plot with a=0.4
\end{eulerprompt}
\eulerimg{17}{Images Use/TugasAplikom4Adi25030630001-031.png}
\begin{eulerprompt}
>plot2d(\{\{"f",0.2\}\},0,1): // plot with a=0.2
\end{eulerprompt}
\eulerimg{17}{Images Use/TugasAplikom4Adi25030630001-032.png}
\begin{eulerprompt}
>plot2d(\{\{"f(x,b)",b=0.1\}\},0,1): // plot with 0.1
\end{eulerprompt}
\eulerimg{17}{Images Use/TugasAplikom4Adi25030630001-033.png}
\begin{eulerprompt}
>function f(x) := x^3-x; ...
>plot2d("f",r=1):
\end{eulerprompt}
\eulerimg{17}{Images Use/TugasAplikom4Adi25030630001-034.png}
\begin{eulercomment}
Berikut adalah rangkuman fungsi yang diterima:

- ekspresi atau ekspresi simbolik dalam variabel x,\\
- fungsi atau fungsi simbolik berdasarkan nama sebagai "f",\\
- fungsi simbolik hanya dengan namanya f.

Fungsi plot2d() juga menerima fungsi simbolik. Untuk fungsi simbolik,
penulisan namanya saja sudah cukup untuk menjalankan fungsi tersebut.
\end{eulercomment}
\begin{eulerprompt}
>function f(x) &= diff(x^x,x)
\end{eulerprompt}
\begin{euleroutput}
  
                              x
                             x  (log(x) + 1)
  
\end{euleroutput}
\begin{eulerprompt}
>plot2d(f,0,2):
\end{eulerprompt}
\eulerimg{17}{Images Use/TugasAplikom4Adi25030630001-035.png}
\begin{eulercomment}
Tentu saja, untuk ekspresi atau ekspresi simbolik, penulisan nama
variabelnya saja sudah cukup untuk menggambar grafiknya.
\end{eulercomment}
\begin{eulerprompt}
>expr &= sin(x)*exp(-x)
\end{eulerprompt}
\begin{euleroutput}
  
                                - x
                               E    sin(x)
  
\end{euleroutput}
\begin{eulerprompt}
>plot2d(expr,0,3pi):
\end{eulerprompt}
\eulerimg{17}{Images Use/TugasAplikom4Adi25030630001-036.png}
\begin{eulerprompt}
>function f(x) &= x^x;
>plot2d(f,r=1,cx=1,cy=1,color=blue,thickness=2);
>plot2d(&diff(f(x),x),>add,color=red,style="-.-"):
\end{eulerprompt}
\eulerimg{17}{Images Use/TugasAplikom4Adi25030630001-037.png}
\begin{eulercomment}
Untuk gaya garis (line style), terdapat beberapa pilihan yang
tersedia.

- style="...": Pilih dari "-", "--", "-.", ".", ".-.", "-.-".\\
- \\
color: Lihat penjelasan warna di bawah ini.

- thickness: Nilai bawaan (default) adalah 1.

Warna dapat dipilih dari salah satu warna standar atau menggunakan
warna RGB.

- 0..15: indeks warna bawaan.\\
- \\
konstanta warna: white, black, red, green, blue, cyan, olive,
lightgray, gray, darkgray, orange, lightgreen, turquoise, lightblue,
lightorange, yellow

- rgb(red,green,blue): parameter berupa angka riil dalam rentang
[0,1].
\end{eulercomment}
\begin{eulerprompt}
>plot2d("exp(-x^2)",r=2,color=red,thickness=3,style="--"):
\end{eulerprompt}
\eulerimg{17}{Images Use/TugasAplikom4Adi25030630001-038.png}
\begin{eulercomment}
Berikut adalah tampilan warna-warna bawaan (*predefined colors*) dari
EMT.
\end{eulercomment}
\begin{eulerprompt}
>aspect(2); columnsplot(ones(1,16),lab=0:15,grid=0,color=0:15):
\end{eulerprompt}
\eulerimg{13}{Images Use/TugasAplikom4Adi25030630001-039.png}
\begin{eulercomment}
Tapi kamu juga bisa memakai warna apa saja yang kamu mau.
\end{eulercomment}
\begin{eulerprompt}
>columnsplot(ones(1,16),grid=0,color=rgb(0,0,linspace(0,1,15))):
\end{eulerprompt}
\eulerimg{13}{Images Use/TugasAplikom4Adi25030630001-040.png}
\eulerheading{Menggambar Beberapa Kurva pada bidang koordinat yang sama}
\begin{eulercomment}
Untuk menggambar lebih dari satu fungsi (banyak fungsi) ke dalam satu
jendela, terdapat beberapa cara yang bisa digunakan. Salah satu
metodenya adalah menggunakan opsi `\textgreater{}add` pada pemanggilan `plot2d`
berikutnya setelah pemanggilan pertama. Fitur ini sebenarnya sudah
kita gunakan pada contoh-contoh di atas.
\end{eulercomment}
\begin{eulerprompt}
>aspect(); plot2d("cos(x)",r=2,grid=6); plot2d("x",style=".",>add):
\end{eulerprompt}
\eulerimg{27}{Images Use/TugasAplikom4Adi25030630001-041.png}
\begin{eulerprompt}
>aspect(1.5); plot2d("sin(x)",0,2pi); plot2d("cos(x)",color=blue,style="--",>add):
\end{eulerprompt}
\eulerimg{17}{Images Use/TugasAplikom4Adi25030630001-042.png}
\begin{eulercomment}
Salah satu kegunaan \textgreater{}add adalah untuk menambahkan titik pada kurva.
\end{eulercomment}
\begin{eulerprompt}
>plot2d("sin(x)",0,pi); plot2d(2,sin(2),>points,>add):
\end{eulerprompt}
\eulerimg{17}{Images Use/TugasAplikom4Adi25030630001-043.png}
\begin{eulercomment}
Kita menambahkan titik potong beserta labelnya (pada posisi "cl" untuk
*center left*), lalu menyisipkan hasilnya ke dalam notebook. Kita juga
menambahkan judul pada grafik tersebut.
\end{eulercomment}
\begin{eulerprompt}
>plot2d(["cos(x)","x"],r=1.1,cx=0.5,cy=0.5, ...
>  color=[black,blue],style=["-","."], ...
>  grid=1);
\end{eulerprompt}
\begin{euleroutput}
  Illegal parameter after named parameter!
  Error in:
  ... ,   color=[black,blue],style=["-","."],   grid=1); ...
                                                       ^
\end{euleroutput}
\begin{eulerprompt}
>x0=solve("cos(x)-x",1);  ...
>  plot2d(x0,x0,>points,>add,title="Intersection Demo");  ...
>  label("cos(x) = x",x0,x0,pos="cl",offset=20):
\end{eulerprompt}
\eulerimg{17}{Images Use/TugasAplikom4Adi25030630001-044.png}
\begin{eulercomment}
Pada demo berikut, kita akan menggambarkan grafik fungsi
sinc(x)=sin(x)/x beserta ekspansi Taylor ke-8 dan ke-16 miliknya. Kita
menghitung ekspansi ini menggunakan Maxima melalui ekspresi simbolik.\\
Grafik ini dibuat dalam perintah multi-baris berikut dengan tiga kali
pemanggilan plot2d(). Pemanggilan kedua dan ketiga memiliki bendera
(flag) \textgreater{}add yang diaktifkan, sehingga grafik tersebut menggunakan
rentang sumbu yang sama seperti sebelumnya.

Kita juga menambahkan kotak label yang menjelaskan fungsi-fungsi
tersebut.

\end{eulercomment}
\begin{eulerprompt}
>$taylor(sin(x)/x,x,0,4)
\end{eulerprompt}
\begin{eulerformula}
\[
\frac{x^4}{120}-\frac{x^2}{6}+1
\]
\end{eulerformula}
\begin{eulerprompt}
>plot2d("sinc(x)",0,4pi,color=green,thickness=2); ...
>  plot2d(&taylor(sin(x)/x,x,0,8),>add,color=blue,style="--"); ...
>  plot2d(&taylor(sin(x)/x,x,0,16),>add,color=red,style="-.-"); ...
>  labelbox(["sinc","T8","T16"],styles=["-","--","-.-"], ...
>    colors=[black,blue,red]):
\end{eulerprompt}
\eulerimg{17}{Images Use/TugasAplikom4Adi25030630001-046.png}
\begin{eulercomment}
Pada contoh berikut, kita menghasilkan polinomial Bernstein (Bernstein
Polynomials).

\end{eulercomment}
\begin{eulerformula}
\[
B_i(x) = \binom{n}{i} x^i (1-x)^{n-i}
\]
\end{eulerformula}
\begin{eulerprompt}
>plot2d("(1-x)^10",0,1); // plot first function
>for i=1 to 10; plot2d("bin(10,i)*x^i*(1-x)^(10-i)",>add); end;
>insimg;
\end{eulerprompt}
\eulerimg{17}{Images Use/TugasAplikom4Adi25030630001-048.png}
\begin{eulercomment}
Metode kedua adalah menggunakan sepasang matriks untuk nilai-x dan
matriks untuk nilai-y yang berukuran sama.

Kita menghasilkan matriks nilai dengan satu Polinomial Bernstein di
setiap barisnya. Untuk ini, kita cukup menggunakan vektor kolom dari
\textdollar{}i\textdollar{}. Silakan lihat bagian pengantar tentang bahasa matriks untuk
mempelajari detail selengkapnya.
\end{eulercomment}
\begin{eulerprompt}
>x=linspace(0,1,500);
>n=10; k=(0:n)'; // n is row vector, k is column vector
>y=bin(n,k)*x^k*(1-x)^(n-k); // y is a matrix then
>plot2d(x,y):
\end{eulerprompt}
\eulerimg{17}{Images Use/TugasAplikom4Adi25030630001-049.png}
\begin{eulercomment}
Perhatikan bahwa parameter warna (*color*) dapat berupa sebuah vektor.
Jika menggunakan vektor, maka setiap warna akan diterapkan pada
masing-masing baris dari matriks tersebut.
\end{eulercomment}
\begin{eulerprompt}
>x=linspace(0,1,200); y=x^(1:10)'; plot2d(x,y,color=1:10):
\end{eulerprompt}
\eulerimg{17}{Images Use/TugasAplikom4Adi25030630001-050.png}
\begin{eulercomment}
Metode lainnya adalah menggunakan vektor dari ekspresi (string). Kamu
kemudian dapat menggunakan matriks/vektor warna, matriks gaya garis
(*styles*), dan matriks ketebalan (*thicknesses*) dengan panjang yang
sama.
\end{eulercomment}
\begin{eulerprompt}
>plot2d(["sin(x)","cos(x)"],0,2pi,color=4:5): 
\end{eulerprompt}
\eulerimg{17}{Images Use/TugasAplikom4Adi25030630001-051.png}
\begin{eulerprompt}
>plot2d(["sin(x)","cos(x)"],0,2pi): // plot vector of expressions
\end{eulerprompt}
\eulerimg{17}{Images Use/TugasAplikom4Adi25030630001-052.png}
\begin{eulercomment}
Kita dapat memperoleh vektor tersebut dari Maxima menggunakan
`makelist()` dan `mxm2str()`.
\end{eulercomment}
\begin{eulerprompt}
>v &= makelist(binomial(10,i)*x^i*(1-x)^(10-i),i,0,10) // make list
\end{eulerprompt}
\begin{euleroutput}
  
                 10            9              8  2             7  3
         [(1 - x)  , 10 (1 - x)  x, 45 (1 - x)  x , 120 (1 - x)  x , 
             6  4             5  5             4  6             3  7
  210 (1 - x)  x , 252 (1 - x)  x , 210 (1 - x)  x , 120 (1 - x)  x , 
            2  8              9   10
  45 (1 - x)  x , 10 (1 - x) x , x  ]
  
\end{euleroutput}
\begin{eulerprompt}
>mxm2str(v) // get a vector of strings from the symbolic vector
\end{eulerprompt}
\begin{euleroutput}
  (1-x)^10
  10*(1-x)^9*x
  45*(1-x)^8*x^2
  120*(1-x)^7*x^3
  210*(1-x)^6*x^4
  252*(1-x)^5*x^5
  210*(1-x)^4*x^6
  120*(1-x)^3*x^7
  45*(1-x)^2*x^8
  10*(1-x)*x^9
  x^10
\end{euleroutput}
\begin{eulerprompt}
>plot2d(mxm2str(v),0,1): // plot functions
\end{eulerprompt}
\eulerimg{17}{Images Use/TugasAplikom4Adi25030630001-053.png}
\begin{eulercomment}
Alternatif lainnya adalah menggunakan bahasa matriks dari Euler.

Jika suatu ekspresi menghasilkan matriks fungsi, dengan satu fungsi di
setiap barisnya, maka semua fungsi ini akan digambarkan ke dalam satu
grafik.

Untuk melakukan hal ini, gunakan vektor parameter dalam bentuk vektor
kolom. Jika matriks/vektor warna ditambahkan, warna tersebut akan
diterapkan pada masing-masing baris grafik.
\end{eulercomment}
\begin{eulerprompt}
>n=(1:10)'; plot2d("x^n",0,1,color=1:10):
\end{eulerprompt}
\eulerimg{17}{Images Use/TugasAplikom4Adi25030630001-054.png}
\begin{eulercomment}
Ekspresi dan fungsi satu baris (one-line functions) dapat membaca
variabel global.

Jika kamu tidak dapat menggunakan variabel global, kamu perlu
menggunakan fungsi dengan parameter tambahan, lalu meneruskan
parameter ini sebagai parameter titik koma (semicolon parameter).

Berhati-hatilah untuk meletakkan semua parameter yang ditugaskan
(assigned parameters) di bagian akhir perintah plot2d. Pada contoh
ini, kita meneruskan a=5 ke fungsi f, yang kita gambarkan dari -10
hingga 10.
\end{eulercomment}
\begin{eulerprompt}
>function f(x,a) := 1/a*exp(-x^2/a); ...
>plot2d("f",-10,10;5,thickness=2,title="a=5"):
\end{eulerprompt}
\eulerimg{17}{Images Use/TugasAplikom4Adi25030630001-055.png}
\begin{eulercomment}
Alternatif lainnya adalah menggunakan *collection* yang berisi nama
fungsi beserta semua parameter tambahannya. Daftar khusus ini disebut
*call collections*, dan ini merupakan cara yang disarankan untuk
meneruskan argumen ke suatu fungsi yang dirinya sendiri diteruskan
sebagai argumen ke fungsi lain.

Pada contoh berikut, kita menggunakan perulangan (*loop*) untuk
menggambar beberapa fungsi sekaligus (silakan lihat tutorial tentang
pemograman untuk penjelasan mengenai perulangan `for`).
\end{eulercomment}
\begin{eulerprompt}
>plot2d(\{\{"f",1\}\},-10,10); ...
>for a=2:10; plot2d(\{\{"f",a\}\},>add); end:
\end{eulerprompt}
\eulerimg{17}{Images Use/TugasAplikom4Adi25030630001-056.png}
\begin{eulercomment}
Kita bisa mencapai hasil yang sama dengan cara berikut menggunakan
bahasa matriks dari EMT. Setiap baris dari matriks `f(x,a)` mewakili
satu fungsi. Selain itu, kita dapat mengatur warna untuk setiap baris
matriks tersebut. Klik dua kali pada fungsi `getspectral()` untuk
melihat penjelasannya.

Hasil double click :\\
figure\\
\end{eulercomment}
\eulersubheading{}
\begin{eulercomment}
getspectral

function getspectral (x)

Function in file : C:\textbackslash{}Program Files\textbackslash{}Euler x64\textbackslash{}util\textbackslash{}statistics

Get a spectral color for 0\textless{}=x\textless{}=1.

The scheme runs from blue (0) to red (1)
\end{eulercomment}
\begin{eulerprompt}
>x=-10:0.01:10; a=(1:10)'; plot2d(x,f(x,a),color=getspectral(a/10)):
\end{eulerprompt}
\eulerimg{17}{Images Use/TugasAplikom4Adi25030630001-057.png}
\eulersubheading{Text Labels}
\begin{eulercomment}
Dekorasi sederhana yang dapat digunakan antara lain:

- judul dengan title="..."\\
- label sumbu-x dan sumbu-y dengan xl="...", yl="..."\\
- label teks tambahan dengan label("...",x,y)

Perintah label akan menggambarkan teks ke dalam grafik yang sedang
aktif pada koordinat grafik \textdollar{}(x,y)\textdollar{}. Perintah ini juga dapat menerima
argumen posisi (positional argument).
\end{eulercomment}
\begin{eulerprompt}
>plot2d("x^3-x",-1,2,title="y=x^3-x",yl="y",xl="x"):
\end{eulerprompt}
\eulerimg{17}{Images Use/TugasAplikom4Adi25030630001-058.png}
\begin{eulerprompt}
>expr := "log(x)/x"; ...
>  plot2d(expr,0.5,5,title="y="+expr,xl="x",yl="y"); ...
>  label("(1,0)",1,0); label("Max",E,expr(E),pos="lc"):
\end{eulerprompt}
\eulerimg{17}{Images Use/TugasAplikom4Adi25030630001-059.png}
\begin{eulercomment}
Ada juga fungsi `labelbox()`, yang dapat menampilkan daftar fungsi
beserta teks penjelasannya. Fungsi ini menerima vektor dari string dan
warna, dengan satu elemen untuk setiap fungsi.
\end{eulercomment}
\begin{eulerprompt}
>function f(x) &= x^2*exp(-x^2);  ...
>plot2d(&f(x),a=-3,b=3,c=-1,d=1);  ...
>plot2d(&diff(f(x),x),>add,color=blue,style="--"); ...
>labelbox(["function","derivative"],styles=["-","--"], ...
>   colors=[black,blue],w=0.4):
\end{eulerprompt}
\eulerimg{17}{Images Use/TugasAplikom4Adi25030630001-060.png}
\begin{eulercomment}
Secara *default*, kotak label ditautkan (*anchored*) di sudut kanan
atas, tetapi opsi `\textgreater{}left` akan menautkannya di sudut kiri atas. Kamu
juga dapat memindahkannya ke posisi mana saja sesuai keinginan. Posisi
tautan tersebut adalah sudut kanan atas dari kotak label, dan nilainya
berupa pecahan (fraksi) dari ukuran jendela grafik. Lebar kotaknya
akan menyesuaikan secara otomatis.

Untuk grafik berupa titik-titik (*point plots*), kotak label ini juga
dapat digunakan. Cukup tambahkan parameter `\textgreater{}points`, atau gunakan
vektor bendera (*vector of flags*) dengan satu bendera untuk
masing-masing label.

Pada contoh berikut, hanya ada satu fungsi yang digambarkan. Jadi,
kita dapat menggunakan string biasa alih-alih vektor string. Kita juga
mengatur warna teks menjadi hitam untuk contoh ini.
\end{eulercomment}
\begin{eulerprompt}
>n=10; plot2d(0:n,bin(n,0:n),>addpoints); ...
>labelbox("Binomials",styles="[]",>points,x=0.1,y=0.1, ...
>tcolor=black,>left):
\end{eulerprompt}
\eulerimg{17}{Images Use/TugasAplikom4Adi25030630001-061.png}
\begin{eulercomment}
Gaya grafik seperti ini juga tersedia pada fungsi `statplot()`. Sama
seperti pada `plot2d()`, warna dapat diatur untuk setiap baris grafik.
Terdapat pula grafik khusus lainnya yang ditujukan untuk keperluan
statistik (silakan lihat tutorial tentang statistik).
\end{eulercomment}
\begin{eulerprompt}
>statplot(1:10,random(2,10),color=[red,blue]):
\end{eulerprompt}
\eulerimg{17}{Images Use/TugasAplikom4Adi25030630001-062.png}
\begin{eulercomment}
Fitur yang serupa dengan ini adalah fungsi `textbox()`.

Secara *default*, lebar kotaknya mengikuti lebar maksimum dari
baris-baris teks di dalamnya. Namun, lebarnya juga dapat diatur secara
manual oleh pengguna.
\end{eulercomment}
\begin{eulerprompt}
>function f(x) &= exp(-x)*sin(2*pi*x); ...
>plot2d("f(x)",0,2pi); ...
>textbox(latex("\(\backslash\)text\{Example of a damped oscillation\}\(\backslash\) f(x)=e^\{-x\}sin(2\(\backslash\)pi x)"),w=0.85):
\end{eulerprompt}
\eulerimg{17}{Images Use/TugasAplikom4Adi25030630001-063.png}
\begin{eulercomment}
Label teks, judul, kotak label, dan teks lainnya dapat memuat string
Unicode (silakan lihat panduan sintaks EMT untuk informasi lebih
lanjut mengenai string Unicode).
\end{eulercomment}
\begin{eulerprompt}
>plot2d("x^3-x",title=u"x &rarr; x&sup3; - x"):
\end{eulerprompt}
\eulerimg{17}{Images Use/TugasAplikom4Adi25030630001-064.png}
\begin{eulercomment}
Label pada sumbu-x dan sumbu-y dapat diposisikan secara vertikal,
demikian pula dengan sumbu itu sendiri.
\end{eulercomment}
\begin{eulerprompt}
>plot2d("sinc(x)",0,2pi,xl="x",yl=u"x &rarr; sinc(x)",>vertical):
\end{eulerprompt}
\eulerimg{17}{Images Use/TugasAplikom4Adi25030630001-065.png}
\eulersubheading{LaTeX}
\begin{eulercomment}
Kamu juga dapat menggambarkan formula LaTeX jika kamu sudah menginstal
sistem LaTeX. Saya menyarankan untuk menggunakan MiKTeX. Jalur
(*path*) menuju berkas biner "latex" dan "dvipng" harus sudah
ditambahkan ke *system path*, atau kamu harus mengaturnya terlebih
dahulu di menu opsi.

Perlu diingat bahwa pemrosesan (*parsing*) LaTeX membutuhkan waktu
yang agak lambat. Jika kamu ingin menggunakan LaTeX dalam grafik
animasi, sebaiknya panggil fungsi `latex()` sekali saja sebelum
perulangan (*loop*) dan gunakan hasilnya (berupa gambar dalam matriks
RGB).

Pada grafik berikut, kita menggunakan LaTeX untuk label sumbu-x dan
sumbu-y, label teks, kotak label (*label box*), serta judul grafik.
\end{eulercomment}
\begin{eulerprompt}
>plot2d("exp(-x)*sin(x)/x",a=0,b=2pi,c=0,d=1,grid=6,color=blue, ...
>  title=latex("\(\backslash\)text\{Function $\(\backslash\)Phi$\}"), ...
>  xl=latex("\(\backslash\)phi"),yl=latex("\(\backslash\)Phi(\(\backslash\)phi)")); ...
>textbox( ...
>  latex("\(\backslash\)Phi(\(\backslash\)phi) = e^\{-\(\backslash\)phi\} \(\backslash\)frac\{\(\backslash\)sin(\(\backslash\)phi)\}\{\(\backslash\)phi\}"),x=0.8,y=0.5); ...
>label(latex("\(\backslash\)Phi",color=blue),1,0.4):
\end{eulerprompt}
\eulerimg{17}{Images Use/TugasAplikom4Adi25030630001-066.png}
\begin{eulercomment}
Sering kali, kita menginginkan jarak (*spacing*) yang tidak sebanding
(*non-conformal*) dan label teks pada sumbu-x. Kita dapat menggunakan
`xaxis()` dan `yaxis()` seperti yang akan ditunjukkan nanti.

Cara termudah adalah dengan membuat grafik kosong beserta bingkainya
menggunakan `grid=4`, lalu menambahkan garis kisi (*grids*) dengan
`ygrid()` dan `xgrid()`. Pada contoh berikut, kita menggunakan tiga
string LaTeX untuk label pada sumbu-x dengan `xtick()`.
\end{eulercomment}
\begin{eulerprompt}
>plot2d("sinc(x)",0,2pi,grid=4,<ticks); ...
>ygrid(-2:0.5:2,grid=6); ...
>xgrid([0:2]*pi,<ticks,grid=6);  ...
>xtick([0,pi,2pi],["0","\(\backslash\)pi","2\(\backslash\)pi"],>latex):
\end{eulerprompt}
\eulerimg{17}{Images Use/TugasAplikom4Adi25030630001-067.png}
\begin{eulercomment}
Tentu saja, fungsi biasa juga dapat digunakan.
\end{eulercomment}
\begin{eulerprompt}
>function map f(x) ...
\end{eulerprompt}
\begin{eulerudf}
  if x>0 then return x^4
  else return x^2
  endif
  endfunction
\end{eulerudf}
\begin{eulercomment}
Parameter `map` membantu agar fungsi dapat digunakan untuk vektor.
Untuk pembuatan grafik, hal ini sebenarnya tidak terlalu diperlukan.
Namun, untuk menunjukkan bahwa vektorisasi (*vectorization*) itu
berguna, kita menambahkan beberapa titik kunci ke grafik pada x = -1,
x = 0, dan x = 1.

Pada grafik berikut, kita juga memasukkan beberapa kode LaTeX. Kode
ini kita gunakan untuk dua buah label dan sebuah kotak teks (*text
box*). Tentu saja, kamu hanya bisa menggunakan LaTeX jika sistem LaTeX
sudah terinstal dengan benar di komputer kamu.
\end{eulercomment}
\begin{eulerprompt}
>plot2d("f",-1,1,xl="x",yl="f(x)",grid=6);  ...
>plot2d([-1,0,1],f([-1,0,1]),>points,>add); ...
>label(latex("x^3"),0.72,f(0.72)); ...
>label(latex("x^2"),-0.52,f(-0.52),pos="ll"); ...
>textbox( ...
>  latex("f(x)=\(\backslash\)begin\{cases\} x^3 & x>0 \(\backslash\)\(\backslash\) x^2 & x \(\backslash\)le 0\(\backslash\)end\{cases\}"), ...
>  x=0.7,y=0.2):
\end{eulerprompt}
\eulerimg{17}{Images Use/TugasAplikom4Adi25030630001-068.png}
\begin{eulercomment}
\end{eulercomment}
\eulersubheading{User Interaction}
\begin{eulercomment}
Saat menggambarkan fungsi atau ekspresi, parameter \textgreater{}user memungkinkan
pengguna untuk melakukan zoom dan menggeser (shift) tampilan grafik
menggunakan tombol panah (cursor keys) atau tetikus (mouse). Pengguna
dapat:

- melakukan zoom dengan tombol + atau -\\
- menggeser grafik menggunakan tombol panah (cursor keys)\\
- memilih area grafik menggunakan mouse\\
- mengembalikan tampilan ke posisi awal dengan tombol spasi (space)\\
- keluar dari mode interaktif dengan tombol Enter (return)

Tombol spasi akan mengembalikan skala grafik ke jendela grafik (plot
window) semula.

Sedangkan saat menggambarkan data, bendera \textgreater{}user hanya akan menunggu
hingga ada penekanan tombol (key stroke) dari pengguna.
\end{eulercomment}
\begin{eulerprompt}
>plot2d(\{\{"x^3-a*x",a=1\}\},>user,title="Press any key!"):
\end{eulerprompt}
\eulerimg{17}{Images Use/TugasAplikom4Adi25030630001-069.png}
\begin{eulerprompt}
>plot2d("exp(x)*sin(x)",user=true, ...
>  title="+/- or cursor keys (return to exit)"):
\end{eulerprompt}
\eulerimg{17}{Images Use/TugasAplikom4Adi25030630001-070.png}
\begin{eulercomment}
Berikut ini menunjukkan cara interaksi pengguna tingkat lanjut
(silakan lihat tutorial tentang pemograman untuk detail lebih lanjut).

Fungsi bawaan `mousedrag()` bertugas menunggu peristiwa (*events*)
dari tetikus (*mouse*) atau papan ketik (*keyboard*). Fungsi ini akan
melaporkan saat tombol tetikus ditekan (*mouse down*), digeser (*mouse
moved*), dilepas (*mouse up*), serta saat ada tombol papan ketik yang
ditekan. Fungsi `dragpoints()` memanfaatkan fitur ini untuk
memungkinkan pengguna menggeser (*drag*) titik mana pun pada grafik.

Pertama-tama, kita memerlukan sebuah fungsi grafik terlebih dahulu.
Sebagai contoh, kita melakukan interpolasi menggunakan polinomial pada
5 titik. Fungsi ini harus menggambar ke dalam area grafik yang tetap
(*fixed plot area*).
\end{eulercomment}
\begin{eulerprompt}
>function plotf(xp,yp,select) ...
\end{eulerprompt}
\begin{eulerudf}
    d=interp(xp,yp);
    plot2d("interpval(xp,d,x)";d,xp,r=2);
    plot2d(xp,yp,>points,>add);
    if select>0 then
      plot2d(xp[select],yp[select],color=red,>points,>add);
    endif;
    title("Drag one point, or press space or return!");
  endfunction
\end{eulerudf}
\begin{eulercomment}
Perhatikan penggunaan parameter titik koma (`;`) pada `plot2d` (`d`
dan `xp`), yang diteruskan ke evaluasi fungsi `interp()`. Tanpa ini,
kita harus menulis fungsi `plotinterp()` terlebih dahulu yang
mengakses nilai-nilai tersebut secara global.

Sekarang kita akan menghasilkan beberapa nilai acak dan membiarkan
pengguna menggeser titik-titik tersebut.
\end{eulercomment}
\begin{eulerprompt}
>t=-1:0.5:1; dragpoints("plotf",t,random(size(t))-0.5):
\end{eulerprompt}
\eulerimg{17}{Images Use/TugasAplikom4Adi25030630001-071.png}
\begin{eulercomment}
Ada juga fungsi yang dapat menggambarkan fungsi lain berdasarkan
vektor parameter, serta memungkinkan pengguna untuk menyesuaikan
parameter-parameter tersebut secara interaktif.

Pertama-tama, kita memerlukan fungsi grafik (*plot function*) terlebih
dahulu.
\end{eulercomment}
\begin{eulerprompt}
>function plotf([a,b]) := plot2d("exp(a*x)*cos(2pi*b*x)",0,2pi;a,b);
\end{eulerprompt}
\begin{eulercomment}
Selanjutnya, kita membutuhkan nama untuk setiap parameter, nilai awal
(*initial values*), dan matriks berukuran \textdollar{}n \textbackslash{}times 2\textdollar{} yang berisi
rentang nilai (*ranges*), serta secara opsional sebuah baris judul
(*heading line*).

Terdapat penggeser (*sliders*) interaktif yang dapat digunakan oleh
pengguna untuk mengatur nilai. Fungsi `dragvalues()` menyediakan fitur
ini.
\end{eulercomment}
\begin{eulerprompt}
>dragvalues("plotf",["a","b"],[-1,2],[[-2,2];[1,10]], ...
>  heading="Drag these values:",hcolor=black):
\end{eulerprompt}
\eulerimg{17}{Images Use/TugasAplikom4Adi25030630001-072.png}
\begin{eulercomment}
Dimungkinkan untuk membatasi nilai yang digeser agar hanya berupa
bilangan bulat (*integers*). Sebagai contoh, kita menuliskan sebuah
fungsi grafik yang menggambarkan polinomial Taylor berderajat \textdollar{}n\textdollar{} dari
fungsi kosinus.
\end{eulercomment}
\begin{eulerprompt}
>function plotf(n) ...
\end{eulerprompt}
\begin{eulerudf}
  plot2d("cos(x)",0,2pi,>square,grid=6);
  plot2d(&"taylor(cos(x),x,0,@n)",color=blue,>add);
  textbox("Taylor polynomial of degree "+n,0.1,0.02,style="t",>left);
  endfunction
\end{eulerudf}
\begin{eulercomment}
Sekarang kita mengizinkan derajat \textdollar{}n\textdollar{} untuk bervariasi dari 0 hingga
20 dalam 20 tahapan. Hasil dari `dragvalues()` digunakan untuk
menggambar sketsa dengan nilai \textdollar{}n\textdollar{} tersebut, serta menyisipkan hasil
grafiknya ke dalam *notebook*.
\end{eulercomment}
\begin{eulerprompt}
>nd=dragvalues("plotf","degree",2,[0,20],20,y=0.8, ...
>   heading="Drag the value:"); ...
>plotf(nd):
\end{eulerprompt}
\eulerimg{17}{Images Use/TugasAplikom4Adi25030630001-073.png}
\begin{eulercomment}
Berikut adalah demonstrasi sederhana dari penggunaan fungsi tersebut.
Pengguna dapat menggambar di atas jendela grafik (*plot window*), yang
akan meninggalkan jejak berupa titik-titik.
\end{eulercomment}
\begin{eulerprompt}
>function dragtest ...
\end{eulerprompt}
\begin{eulerudf}
    plot2d(none,r=1,title="Drag with the mouse, or press any key!");
    start=0;
    repeat
      \{flag,m,time\}=mousedrag();
      if flag==0 then return; endif;
      if flag==2 then
        hold on; mark(m[1],m[2]); hold off;
      endif;
    end
  endfunction
\end{eulerudf}
\begin{eulerprompt}
>dragtest // lihat hasilnya dan cobalah lakukan!
\end{eulerprompt}
\eulersubheading{Styles of 2D Plots}
\begin{eulercomment}
Secara bawaan (*by default*), EMT akan menghitung tanda sumbu (*axis
ticks*) secara otomatis dan menambahkan label di setiap tandanya.
Pengaturan ini dapat diubah menggunakan parameter `grid`. Gaya standar
dari sumbu dan label juga bisa disesuaikan. Selain itu, kamu bisa
menambahkan label dan judul secara manual. Untuk mengembalikan semua
gaya ke tampilan awal (*default*), gunakan fungsi `reset()`.
\end{eulercomment}
\begin{eulerprompt}
>aspect();
>figure(3,4); ...
> figure(1); plot2d("x^3-x",grid=0); ... // no grid, frame or axis
> figure(2); plot2d("x^3-x",grid=1); ... // x-y-axis
> figure(3); plot2d("x^3-x",grid=2); ... // default ticks
> figure(4); plot2d("x^3-x",grid=3); ... // x-y- axis with labels inside
> figure(5); plot2d("x^3-x",grid=4); ... // no ticks, only labels
> figure(6); plot2d("x^3-x",grid=5); ... // default, but no margin
> figure(7); plot2d("x^3-x",grid=6); ... // axes only
> figure(8); plot2d("x^3-x",grid=7); ... // axes only, ticks at axis
> figure(9); plot2d("x^3-x",grid=8); ... // axes only, finer ticks at axis
> figure(10); plot2d("x^3-x",grid=9); ... // default, small ticks inside
> figure(11); plot2d("x^3-x",grid=10); ...// no ticks, axes only
> figure(0):
\end{eulerprompt}
\eulerimg{27}{Images Use/TugasAplikom4Adi25030630001-074.png}
\begin{eulercomment}
Parameter `\textless{}frame` berfungsi untuk mematikan bingkai (*frame*), dan
`framecolor=blue` akan mengatur warna bingkai menjadi biru.

Jika kamu ingin menentukan tanda sumbu (*ticks*) buatanmu sendiri,
kamu bisa menggunakan `style=0` dan menambahkan semua elemennya secara
manual nanti.
\end{eulercomment}
\begin{eulerprompt}
>aspect(1.5); 
>plot2d("x^3-x",grid=0); // plot
>frame; xgrid([-1,0,1]); ygrid(0): // add frame and grid
\end{eulerprompt}
\eulerimg{17}{Images Use/TugasAplikom4Adi25030630001-075.png}
\begin{eulercomment}
Untuk judul grafik (*plot title*) dan label sumbu (*axis labels*),
silakan lihat contoh berikut.
\end{eulercomment}
\begin{eulerprompt}
>plot2d("exp(x)",-1,1);
>textcolor(black); // set the text color to black
>title(latex("y=e^x")); // title above the plot
>xlabel(latex("x")); // "x" for x-axis
>ylabel(latex("y"),>vertical); // vertical "y" for y-axis
>label(latex("(0,1)"),0,1,color=blue): // label a point
\end{eulerprompt}
\eulerimg{17}{Images Use/TugasAplikom4Adi25030630001-076.png}
\begin{eulercomment}
Sumbu grafik juga dapat digambar secara terpisah menggunakan fungsi
`xaxis()` dan `yaxis()`.
\end{eulercomment}
\begin{eulerprompt}
>plot2d("x^3-x",<grid,<frame);
>xaxis(0,xx=-2:1,style="->"); yaxis(0,yy=-5:5,style="->"):
\end{eulerprompt}
\eulerimg{17}{Images Use/TugasAplikom4Adi25030630001-077.png}
\begin{eulercomment}
Teks pada grafik dapat ditambahkan dengan menggunakan `label()`. Pada
contoh berikut, `"lc"` berarti *lower center* (tengah bawah).
Parameter ini mengatur posisi label relatif terhadap koordinat grafik.
\end{eulercomment}
\begin{eulerprompt}
>function f(x) &= x^3-x
\end{eulerprompt}
\begin{euleroutput}
  
                                   3
                                  x  - x
  
\end{euleroutput}
\begin{eulerprompt}
>plot2d(f,-1,1,>square);
>x0=fmin(f,0,1); // compute point of minimum
>label("Rel. Min.",x0,f(x0),pos="lc"): // add a label there
\end{eulerprompt}
\eulerimg{17}{Images Use/TugasAplikom4Adi25030630001-078.png}
\begin{eulercomment}
Terdapat juga fitur kotak teks (*text boxes*).
\end{eulercomment}
\begin{eulerprompt}
>plot2d(&f(x),-1,1,-2,2); // function
>plot2d(&diff(f(x),x),>add,style="--",color=red); // derivative
>labelbox(["f","f'"],["-","--"],[black,red]): // label box
\end{eulerprompt}
\eulerimg{17}{Images Use/TugasAplikom4Adi25030630001-079.png}
\begin{eulerprompt}
>plot2d(["exp(x)","1+x"],color=[black,blue],style=["-","-.-"]):
\end{eulerprompt}
\eulerimg{17}{Images Use/TugasAplikom4Adi25030630001-080.png}
\begin{eulerprompt}
>gridstyle("->",color=gray,textcolor=gray,framecolor=gray);  ...
> plot2d("x^3-x",grid=1);   ...
> settitle("y=x^3-x",color=black); ...
> label("x",2,0,pos="bc",color=gray);  ...
> label("y",0,6,pos="cl",color=gray); ...
> reset():
\end{eulerprompt}
\eulerimg{27}{Images Use/TugasAplikom4Adi25030630001-081.png}
\begin{eulercomment}
Untuk kendali yang jauh lebih fleksibel, sumbu-x dan sumbu-y dapat
dibuat secara manual.

Perintah `fullwindow()` berfungsi untuk memperluas jendela grafik
(*plot window*) karena kita tidak lagi membutuhkan ruang untuk label
di luar jendela tersebut. Gunakan `shrinkwindow()` atau `reset()`
untuk mengembalikan ukurannya ke pengaturan awal (*default*).
\end{eulercomment}
\begin{eulerprompt}
>fullwindow; ...
> gridstyle(color=darkgray,textcolor=darkgray); ...
> plot2d(["2^x","1","2^(-x)"],a=-2,b=2,c=0,d=4,<grid,color=4:6,<frame); ...
> xaxis(0,-2:1,style="->"); xaxis(0,2,"x",<axis); ...
> yaxis(0,4,"y",style="->"); ...
> yaxis(-2,1:4,>left); ...
> yaxis(2,2^(-2:2),style=".",<left); ...
> labelbox(["2^x","1","2^-x"],colors=4:6,x=0.8,y=0.2); ...
> reset:
\end{eulerprompt}
\eulerimg{27}{Images Use/TugasAplikom4Adi25030630001-082.png}
\begin{eulercomment}
Berikut adalah contoh lainnya, di mana *string* Unicode digunakan
serta sumbu ditempatkan di luar area grafik (*plot area*).
\end{eulercomment}
\begin{eulerprompt}
>aspect(1.5); 
>plot2d(["sin(x)","cos(x)"],0,2pi,color=[red,green],<grid,<frame); ...
> xaxis(-1.1,(0:2)*pi,xt=["0",u"&pi;",u"2&pi;"],style="-",>ticks,>zero);  ...
> xgrid((0:0.5:2)*pi,<ticks); ...
> yaxis(-0.1*pi,-1:0.2:1,style="-",>zero,>grid); ...
> labelbox(["sin","cos"],colors=[red,green],x=0.5,y=0.2,>left); ...
> xlabel(u"&phi;"); ylabel(u"f(&phi;)"):
\end{eulerprompt}
\eulerimg{17}{Images Use/TugasAplikom4Adi25030630001-083.png}
\eulerheading{Plotting 2D Data}
\begin{eulercomment}
Jika x dan y adalah vektor data, data tersebut akan digunakan sebagai
koordinat-x dan koordinat-y dari sebuah kurva. Dalam hal ini,
parameter a, b, c, dan d, atau jari-jari r dapat ditentukan, atau
jendela grafik (*plot window*) akan menyesuaikan ukurannya secara
otomatis mengikuti data tersebut. Opsi lainnya, `\textgreater{}square` dapat diatur
untuk mempertahankan rasio aspek persegi (*square aspect ratio*).

Menggambar grafik dari suatu ekspresi sebenarnya hanyalah bentuk
singkat dari grafik data (*data plots*). Untuk grafik data, kamu
memerlukan satu atau beberapa baris nilai-x, serta satu atau beberapa
baris nilai-y. Dari rentang (*range*) dan nilai-x yang ada, fungsi
`plot2d` akan menghitung data yang akan digambar, secara bawaan
menggunakan evaluasi adaptif dari fungsi tersebut. Untuk grafik titik
(*point plots*), gunakan `\textgreater{}points`, sedangkan untuk gabungan garis dan
titik, gunakan `\textgreater{}addpoints`.

Namun, kamu juga bisa memasukkan data secara langsung.

- Gunakan vektor baris untuk x dan y untuk satu fungsi.\\
- Matriks untuk x dan y akan digambar baris demi baris.

Berikut adalah contoh dengan satu baris untuk x dan y.
\end{eulercomment}
\begin{eulerprompt}
>x=-10:0.1:10; y=exp(-x^2)*x; plot2d(x,y):
\end{eulerprompt}
\eulerimg{17}{Images Use/TugasAplikom4Adi25030630001-084.png}
\begin{eulercomment}
Data juga dapat digambarkan sebagai titik-titik (*points*). Gunakan
`points=true` untuk opsi ini. Grafik titik ini bekerja layaknya
poligon, namun hanya menggambar bagian sudut-sudutnya saja.

- `style="..."`: Pilih simbol dari opsi yang tersedia: `"[]"`, `"\textless{}\textgreater{}"`,
`"o"`, `"."`, `".."`", `"+"`, ` "*"`, `"[]#"`, `"\textless{}\textgreater{}#"`, `"o#"`,
`"..#"`, `"#"`, `"\textbar{}"`.

Untuk menggambarkan kumpulan titik, gunakan opsi `\textgreater{}points`. Jika
parameter warna berupa vektor warna, maka setiap titik akan
mendapatkan warna yang berbeda. Untuk matriks koordinat dan vektor
kolom warna, warna tersebut akan diterapkan pada setiap baris matriks.

Parameter `\textgreater{}addpoints` berfungsi untuk menambahkan titik-titik pada
segmen garis untuk grafik data.
\end{eulercomment}
\begin{eulerprompt}
>xdata=[1,1.5,2.5,3,4]; ydata=[3,3.1,2.8,2.9,2.7]; // data
>plot2d(xdata,ydata,a=0.5,b=4.5,c=2.5,d=3.5,style="."); // lines
>plot2d(xdata,ydata,>points,>add,style="o"): // add points
\end{eulerprompt}
\eulerimg{17}{Images Use/TugasAplikom4Adi25030630001-085.png}
\begin{eulerprompt}
>p=polyfit(xdata,ydata,1); // get regression line
>plot2d("polyval(p,x)",>add,color=red): // add plot of line
\end{eulerprompt}
\eulerimg{17}{Images Use/TugasAplikom4Adi25030630001-086.png}
\eulerheading{Menggambar Daerah Yang Dibatasi Kurva}
\begin{eulercomment}
Grafik data sebenarnya adalah poligon. Kita juga dapat membuat grafik
kurva (*curves*) atau kurva terisi (*filled curves*).

- `filled=true`: Mengisi area di dalam grafik.\\
- `style="..."`: Pilih gaya arsiran dari opsi `"#"`, `"/"`, `"\textbackslash{}"`,
atau `"\textbackslash{}\textbackslash{}/"`.\\
- `fillcolor`: Lihat opsi warna yang tersedia di atas untuk menentukan
warna isian.

Warna isian ditentukan oleh argumen `fillcolor`, dan opsi `\textless{}outline`
secara opsional mencegah penggambaran garis batas (*boundary*) untuk
semua gaya, kecuali gaya bawaan (*default*).
\end{eulercomment}
\begin{eulerprompt}
>t=linspace(0,2pi,1000); // parameter for curve
>x=sin(t)*exp(t/pi); y=cos(t)*exp(t/pi); // x(t) and y(t)
>figure(1,2); aspect(16/9)
>figure(1); plot2d(x,y,r=10); // plot curve
>figure(2); plot2d(x,y,r=10,>filled,style="/",fillcolor=red); // fill curve
>figure(0):
\end{eulerprompt}
\eulerimg{14}{Images Use/TugasAplikom4Adi25030630001-087.png}
\begin{eulercomment}
Pada contoh-contoh berikut, kita menggambar sebuah elips terisi
(*filled ellipse*) dan dua segi enam terisi (*filled hexagons*)
menggunakan kurva tertutup dengan 6 titik serta gaya isian (*fill
style*) yang berbeda.
\end{eulercomment}
\begin{eulerprompt}
>x=linspace(0,2pi,1000); plot2d(sin(x),cos(x)*0.5,r=1,>filled,style="/"):
\end{eulerprompt}
\eulerimg{14}{Images Use/TugasAplikom4Adi25030630001-088.png}
\begin{eulerprompt}
>t=linspace(0,2pi,6); ...
>plot2d(cos(t),sin(t),>filled,style="/",fillcolor=red,r=1.2):
\end{eulerprompt}
\eulerimg{14}{Images Use/TugasAplikom4Adi25030630001-089.png}
\begin{eulerprompt}
>t=linspace(0,2pi,6); plot2d(cos(t),sin(t),>filled,style="#"):
\end{eulerprompt}
\eulerimg{14}{Images Use/TugasAplikom4Adi25030630001-090.png}
\begin{eulercomment}
Contoh lainnya adalah segi tujuh (*septagon*), yang kita buat
menggunakan 7 titik pada lingkaran satuan (*unit circle*).
\end{eulercomment}
\begin{eulerprompt}
>t=linspace(0,2pi,7);  ...
> plot2d(cos(t),sin(t),r=1,>filled,style="/",fillcolor=red):
\end{eulerprompt}
\eulerimg{14}{Images Use/TugasAplikom4Adi25030630001-091.png}
\begin{eulercomment}
Berikut adalah himpunan nilai maksimum dari empat kondisi linear yang
kurang dari atau sama dengan 3. Ini dinyatakan sebagai \textdollar{}A[k] \textbackslash{}cdot v
\textbackslash{}le 3\textdollar{} untuk setiap baris dari \textdollar{}A\textdollar{}. Untuk mendapatkan sudut yang halus
(*nice corners*), kita menggunakan \textdollar{}n\textdollar{} yang relatif besar.
\end{eulercomment}
\begin{eulerprompt}
>A=[2,1;1,2;-1,0;0,-1];
>function f(x,y) := max([x,y].A');
>plot2d("f",r=4,level=[0;3],color=green,n=111):
\end{eulerprompt}
\eulerimg{14}{Images Use/TugasAplikom4Adi25030630001-092.png}
\begin{eulercomment}
Poin utama dari bahasa matriks adalah bahwa ia memungkinkan kita untuk
membuat tabel fungsi secara mudah.
\end{eulercomment}
\begin{eulerprompt}
>t=linspace(0,2pi,1000); x=cos(3*t); y=sin(4*t);
\end{eulerprompt}
\begin{eulercomment}
Sekarang kita memiliki vektor nilai \textdollar{}x\textdollar{} dan \textdollar{}y\textdollar{}. Fungsi `plot2d()`
dapat menggambarkan nilai-nilai ini sebagai kurva yang menghubungkan
titik-titiknya. Grafik tersebut juga dapat diisi (*filled*). Dalam hal
ini, opsi tersebut menghasilkan tampilan yang menarik berkat aturan
winding (*winding rule*) yang digunakan untuk pengisiannya.
\end{eulercomment}
\begin{eulerprompt}
>plot2d(x,y,<grid,<frame,>filled):
\end{eulerprompt}
\eulerimg{14}{Images Use/TugasAplikom4Adi25030630001-093.png}
\begin{eulercomment}
Vektor interval digambarkan terhadap nilai-nilai \textdollar{}x\textdollar{} sebagai wilayah
yang terisi (*filled region*) di antara nilai bawah (*lower values*)
dan nilai atas (*upper values*) dari interval tersebut.

Fitur ini sangat berguna untuk menggambarkan kesalahan atau margin
galat (*errors*) dari hasil perhitungan. Selain itu, opsi ini juga
dapat digunakan untuk menampilkan galat statistik (*statistical
errors*).
\end{eulercomment}
\begin{eulerprompt}
>t=0:0.1:1; ...
> plot2d(t,interval(t-random(size(t)),t+random(size(t))),style="|");  ...
> plot2d(t,t,add=true):
\end{eulerprompt}
\eulerimg{14}{Images Use/TugasAplikom4Adi25030630001-094.png}
\begin{eulercomment}
Jika x adalah vektor yang terurut, dan \textdollar{}y\textdollar{} adalah vektor interval,
maka `plot2d` akan menggambarkan area interval yang terisi (*filled
ranges*) pada bidang grafik. Gaya isiannya (*fill styles*) sama dengan
gaya yang digunakan pada poligon.
\end{eulercomment}
\begin{eulerprompt}
>t=-1:0.01:1; x=~t-0.01,t+0.01~; y=x^3-x;
>plot2d(t,y):
\end{eulerprompt}
\eulerimg{14}{Images Use/TugasAplikom4Adi25030630001-095.png}
\begin{eulercomment}
Memungkinkan juga untuk mengisi area (*regions*) nilai bagi suatu
fungsi tertentu. Untuk melakukan hal ini, parameter `level` harus
berupa matriks berukuran \textdollar{}2 \textbackslash{}times n\textdollar{}. Baris pertama berisi
batas-batas bawah (*lower bounds*), sedangkan baris kedua berisi
batas-batas atas (*upper bounds*).
\end{eulercomment}
\begin{eulerprompt}
>expr := "2*x^2+x*y+3*y^4+y"; // define an expression f(x,y)
>plot2d(expr,level=[0;1],style="-",color=blue): // 0 <= f(x,y) <= 1
\end{eulerprompt}
\eulerimg{14}{Images Use/TugasAplikom4Adi25030630001-096.png}
\begin{eulercomment}
Kita juga bisa mengisi rentang nilai seperti


\end{eulercomment}
\begin{eulerformula}
\[
-1 \le (x^2+y^2)^2-x^2+y^2 \le 0.
\]
\end{eulerformula}
\begin{eulercomment}
\end{eulercomment}
\begin{eulerprompt}
>plot2d("(x^2+y^2)^2-x^2+y^2",r=1.2,level=[-1;0],style="/"):
\end{eulerprompt}
\eulerimg{14}{Images Use/TugasAplikom4Adi25030630001-098.png}
\begin{eulerprompt}
>plot2d("cos(x)","sin(x)^3",xmin=0,xmax=2pi,>filled,style="/"):
\end{eulerprompt}
\eulerimg{14}{Images Use/TugasAplikom4Adi25030630001-099.png}
\eulerheading{Grafik Fungsi Parametrik}
\begin{eulercomment}
Nilai-\textdollar{}x\textdollar{} tidak harus terurut. Pasangan \textdollar{}(x,y)\textdollar{} hanya menggambarkan
sebuah kurva. Jika nilai \textdollar{}x\textdollar{} terurut, maka kurva tersebut merupakan
grafik dari suatu fungsi.

Pada contoh berikut, kita menggambarkan spiral

\end{eulercomment}
\begin{eulerformula}
\[
\gamma(t) = t \cdot (\cos(2\pi t),\sin(2\pi t))
\]
\end{eulerformula}
\begin{eulercomment}
Kita dapat menggunakan titik yang sangat banyak agar tampilannya
halus, atau menggunakan fungsi `adaptive()` untuk mengevaluasi
ekspresi tersebut (lihat fungsi `adaptive()` untuk detail lebih
lanjut).
\end{eulercomment}
\begin{eulerprompt}
>t=linspace(0,1,1000); ...
>plot2d(t*cos(2*pi*t),t*sin(2*pi*t),r=1):
\end{eulerprompt}
\eulerimg{14}{Images Use/TugasAplikom4Adi25030630001-101.png}
\begin{eulercomment}
Sebagai alternatif, kamu juga dapat menggunakan dua ekspresi untuk
membuat kurva. Contoh berikut menggambarkan kurva yang sama seperti di
atas.
\end{eulercomment}
\begin{eulerprompt}
>plot2d("x*cos(2*pi*x)","x*sin(2*pi*x)",xmin=0,xmax=1,r=1):
\end{eulerprompt}
\eulerimg{14}{Images Use/TugasAplikom4Adi25030630001-102.png}
\begin{eulerprompt}
>t=linspace(0,1,1000); r=exp(-t); x=r*cos(2pi*t); y=r*sin(2pi*t);
>plot2d(x,y,r=1):
\end{eulerprompt}
\eulerimg{14}{Images Use/TugasAplikom4Adi25030630001-103.png}
\begin{eulercomment}
Pada contoh berikutnya, kita menggambarkan kurva

\end{eulercomment}
\begin{eulerformula}
\[
\gamma(t) = (r(t) \cos(t), r(t) \sin(t))
\]
\end{eulerformula}
\begin{eulercomment}
dengan

\end{eulercomment}
\begin{eulerformula}
\[
r(t) = 1 + \dfrac{\sin(3t)}{2}.
\]
\end{eulerformula}
\begin{eulerprompt}
>t=linspace(0,2pi,1000); r=1+sin(3*t)/2; x=r*cos(t); y=r*sin(t); ...
>plot2d(x,y,>filled,fillcolor=red,style="/",r=1.5):
\end{eulerprompt}
\eulerimg{14}{Images Use/TugasAplikom4Adi25030630001-106.png}
\eulerheading{Menggambar Grafik Bilangan Kompleks}
\begin{eulercomment}
Larik (*array*) angka kompleks juga dapat digambarkan. Dengan
demikian, titik-titik kisi (*grid points*) akan saling dihubungkan.
Jika jumlah garis kisi ditentukan (atau menggunakan vektor \textdollar{}1 \textbackslash{}times
2\textdollar{} untuk garis kisi) pada argumen `cgrid`, hanya garis-garis kisi
tersebut yang akan ditampilkan.

Matriks angka kompleks akan digambarkan secara otomatis sebagai kisi
pada bidang kompleks (*complex plane*).

Pada contoh berikut, kita menggambarkan bayangan dari lingkaran satuan
(*unit circle*) di bawah fungsi eksponensial. Parameter `cgrid`
menyembunyikan beberapa kurva kisi tersebut.
\end{eulercomment}
\begin{eulerprompt}
>aspect(); r=linspace(0,1,50); a=linspace(0,2pi,80)'; z=r*exp(I*a);...
>plot2d(z,a=-1.25,b=1.25,c=-1.25,d=1.25,cgrid=10):
\end{eulerprompt}
\eulerimg{27}{Images Use/TugasAplikom4Adi25030630001-107.png}
\begin{eulerprompt}
>aspect(1.25); r=linspace(0,1,50); a=linspace(0,2pi,200)'; z=r*exp(I*a);
>plot2d(exp(z),cgrid=[40,10]):
\end{eulerprompt}
\eulerimg{21}{Images Use/TugasAplikom4Adi25030630001-108.png}
\begin{eulerprompt}
>r=linspace(0,1,10); a=linspace(0,2pi,40)'; z=r*exp(I*a);
>plot2d(exp(z),>points,>add):
\end{eulerprompt}
\eulerimg{21}{Images Use/TugasAplikom4Adi25030630001-109.png}
\begin{eulercomment}
Vektor bilangan kompleks digambarkan secara otomatis sebagai kurva
pada bidang kompleks (complex plane) dengan bagian real dan bagian
imajiner.

Pada contoh ini, kita menggambar lingkaran satuan (unit circle) dengan

\end{eulercomment}
\begin{eulerformula}
\[
\gamma(t) = e^{it}
\]
\end{eulerformula}
\begin{eulerprompt}
>t=linspace(0,2pi,1000); ...
>plot2d(exp(I*t)+exp(4*I*t),r=2):
\end{eulerprompt}
\eulerimg{21}{Images Use/TugasAplikom4Adi25030630001-111.png}
\eulerheading{Statistical Plots}
\begin{eulercomment}
Terdapat banyak fungsi yang dikhususkan untuk grafik statistik. Salah
satu grafik yang paling sering digunakan adalah grafik kolom (*column
plot*).

Jumlah kumulatif (*cumulative sum*) dari nilai-nilai terdistribusi
normal (0,1) menghasilkan langkah acak (*random walk*).
\end{eulercomment}
\begin{eulerprompt}
>plot2d(cumsum(randnormal(1,1000))):
\end{eulerprompt}
\eulerimg{21}{Images Use/TugasAplikom4Adi25030630001-112.png}
\begin{eulercomment}
Menggunakan dua baris akan menampilkan langkah acak (*random walk*)
dalam dua dimensi.
\end{eulercomment}
\begin{eulerprompt}
>X=cumsum(randnormal(2,1000)); plot2d(X[1],X[2]):
\end{eulerprompt}
\eulerimg{21}{Images Use/TugasAplikom4Adi25030630001-113.png}
\begin{eulerprompt}
>columnsplot(cumsum(random(10)),style="/",color=blue):
\end{eulerprompt}
\eulerimg{21}{Images Use/TugasAplikom4Adi25030630001-114.png}
\begin{eulercomment}
Opsi ini juga dapat menampilkan string teks sebagai label (*labels*).
\end{eulercomment}
\begin{eulerprompt}
>months=["Jan","Feb","Mar","Apr","May","Jun", ...
>  "Jul","Aug","Sep","Oct","Nov","Dec"];
>values=[10,12,12,18,22,28,30,26,22,18,12,8];
>columnsplot(values,lab=months,color=red,style="-");
>title("Temperature"):
\end{eulerprompt}
\eulerimg{21}{Images Use/TugasAplikom4Adi25030630001-115.png}
\begin{eulerprompt}
>k=0:10;
>plot2d(k,bin(10,k),>bar):
\end{eulerprompt}
\eulerimg{21}{Images Use/TugasAplikom4Adi25030630001-116.png}
\begin{eulerprompt}
>plot2d(k,bin(10,k)); plot2d(k,bin(10,k),>points,>add):
\end{eulerprompt}
\eulerimg{21}{Images Use/TugasAplikom4Adi25030630001-117.png}
\begin{eulerprompt}
>plot2d(normal(1000),normal(1000),>points,grid=6,style=".."):
\end{eulerprompt}
\eulerimg{21}{Images Use/TugasAplikom4Adi25030630001-118.png}
\begin{eulerprompt}
>plot2d(normal(1,1000),>distribution,style="O"):
\end{eulerprompt}
\eulerimg{21}{Images Use/TugasAplikom4Adi25030630001-119.png}
\begin{eulerprompt}
>plot2d("qnormal",0,5;2.5,0.5,>filled):
\end{eulerprompt}
\eulerimg{21}{Images Use/TugasAplikom4Adi25030630001-120.png}
\begin{eulercomment}
Untuk menggambarkan distribusi statistik eksperimental, kamu dapat
menggunakan argumen `distribution=n` pada `plot2d()`.
\end{eulercomment}
\begin{eulerprompt}
>w=randexponential(1,1000); // exponential distribution
>plot2d(w,>distribution): // or distribution=n with n intervals
\end{eulerprompt}
\eulerimg{21}{Images Use/TugasAplikom4Adi25030630001-121.png}
\begin{eulercomment}
Atau kamu dapat menghitung distribusi dari data tersebut dan
menggambar hasilnya menggunakan argumen `\textgreater{}bar` pada `plot3d`, atau
dengan grafik kolom (*column plot*).
\end{eulercomment}
\begin{eulerprompt}
>w=normal(1000); // 0-1-normal distribution
>\{x,y\}=histo(w,10,v=[-6,-4,-2,-1,0,1,2,4,6]); // interval bounds v
>plot2d(x,y,>bar):
\end{eulerprompt}
\eulerimg{21}{Images Use/TugasAplikom4Adi25030630001-122.png}
\begin{eulercomment}
Fungsi `statplot()` mengatur gaya grafik (*style*) menggunakan string
sederhana.
\end{eulercomment}
\begin{eulerprompt}
>statplot(1:10,cumsum(random(10)),"b"):
\end{eulerprompt}
\eulerimg{21}{Images Use/TugasAplikom4Adi25030630001-123.png}
\begin{eulerprompt}
>n=10; i=0:n; ...
>plot2d(i,bin(n,i)/2^n,a=0,b=10,c=0,d=0.3); ...
>plot2d(i,bin(n,i)/2^n,points=true,style="ow",add=true,color=blue):
\end{eulerprompt}
\eulerimg{21}{Images Use/TugasAplikom4Adi25030630001-124.png}
\begin{eulercomment}
Selain itu, data juga dapat digambarkan dalam bentuk batang (*bars*).
Dalam hal ini, \textdollar{}x\textdollar{} harus terurut dan memiliki jumlah elemen satu lebih
banyak daripada \textdollar{}y\textdollar{}. Batang-batang tersebut akan membentang dari
\textdollar{}x[i]\textdollar{} hingga \textdollar{}x[i+1]\textdollar{} dengan nilai \textdollar{}y[i]\textdollar{}. Jika \textdollar{}x\textdollar{} memiliki ukuran
yang sama dengan \textdollar{}y\textdollar{}, maka \textdollar{}x\textdollar{} akan diperpanjang satu elemen secara
otomatis dengan jarak (*spacing*) yang sama seperti elemen terakhir.

Gaya isian (*fill styles*) dapat digunakan persis seperti pada contoh
di atas.
\end{eulercomment}
\begin{eulerprompt}
>n=10; k=bin(n,0:n); ...
>plot2d(-0.5:n+0.5,k,bar=true,fillcolor=lightgray):
\end{eulerprompt}
\eulerimg{21}{Images Use/TugasAplikom4Adi25030630001-125.png}
\begin{eulercomment}
Data untuk grafik batang (`bar=1`) dan histogram (`histogram=1`) dapat
diberikan secara eksplisit melalui `xv` dan `yv`, atau dihitung
langsung dari distribusi eksperimental pada `xv` menggunakan argumen
`\textgreater{}distribution` (atau `distribution=n`). Histogram dari nilai-nilai
`xv` akan dihitung secara otomatis jika menggunakan `\textgreater{}histogram`. Jika
argumen `\textgreater{}even` ditentukan, maka nilai-nilai `xv` akan dihitung dalam
rentang interval bilangan bulat (*integer intervals*).
\end{eulercomment}
\begin{eulerprompt}
>plot2d(normal(10000),distribution=50):
\end{eulerprompt}
\eulerimg{21}{Images Use/TugasAplikom4Adi25030630001-126.png}
\begin{eulerprompt}
>k=0:10; m=bin(10,k); x=(0:11)-0.5; plot2d(x,m,>bar):
\end{eulerprompt}
\eulerimg{21}{Images Use/TugasAplikom4Adi25030630001-127.png}
\begin{eulerprompt}
>columnsplot(m,k):
\end{eulerprompt}
\eulerimg{21}{Images Use/TugasAplikom4Adi25030630001-128.png}
\begin{eulerprompt}
>plot2d(random(600)*6,histogram=6):
\end{eulerprompt}
\eulerimg{21}{Images Use/TugasAplikom4Adi25030630001-129.png}
\begin{eulercomment}
Untuk distribusi data, terdapat parameter `distribution=n` yang akan
menghitung nilai-nilai secara otomatis dan menampilkan distribusi
relatifnya dengan \textdollar{}n\textdollar{} sub-interval.
\end{eulercomment}
\begin{eulerprompt}
>plot2d(normal(1,1000),distribution=10,style="\(\backslash\)/"):
\end{eulerprompt}
\eulerimg{21}{Images Use/TugasAplikom4Adi25030630001-130.png}
\begin{eulercomment}
Dengan parameter `even=true`, ini akan menggunakan interval bilangan
bulat (*integer intervals*).
\end{eulercomment}
\begin{eulerprompt}
>plot2d(intrandom(1,1000,10),distribution=10,even=true):
\end{eulerprompt}
\eulerimg{21}{Images Use/TugasAplikom4Adi25030630001-131.png}
\begin{eulercomment}
Perlu diingat bahwa terdapat banyak grafik statistik lainnya yang
sangat berguna. Kamu bisa melihat panduan/tutorial tentang statistik
untuk pembahasan yang lebih lengkap.
\end{eulercomment}
\begin{eulerprompt}
>columnsplot(getmultiplicities(1:6,intrandom(1,6000,6))):
\end{eulerprompt}
\eulerimg{21}{Images Use/TugasAplikom4Adi25030630001-132.png}
\begin{eulerprompt}
>plot2d(normal(1,1000),>distribution); ...
>  plot2d("qnormal(x)",color=red,thickness=2,>add):
\end{eulerprompt}
\eulerimg{21}{Images Use/TugasAplikom4Adi25030630001-133.png}
\begin{eulercomment}
Terdapat juga berbagai grafik khusus untuk statistik. *Boxplot*
menampilkan kuartil dari distribusi data serta pencilan (*outliers*).
Secara definisi, pencilan pada *boxplot* adalah data yang nilainya
melebihi 1,5 kali rentang 50\% data bagian tengah (*interquartile
range*).
\end{eulercomment}
\begin{eulerprompt}
>M=normal(5,1000); boxplot(quartiles(M)):
\end{eulerprompt}
\eulerimg{21}{Images Use/TugasAplikom4Adi25030630001-134.png}
\eulerheading{Implicit Functions}
\begin{eulercomment}
Grafik implisit (*implicit plots*) menampilkan garis kontur/tingkat
(*level lines*) yang menyelesaikan persamaan \textdollar{}f(x,y)=\textbackslash{}text\{level\}\textdollar{}, di
mana `level` dapat berupa nilai tunggal atau vektor nilai. Jika
`level="auto"`, akan terdapat `nc` garis kontur yang tersebar secara
merata antara nilai minimum dan maksimum dari fungsi tersebut. Warna
yang lebih gelap atau lebih terang dapat ditambahkan dengan argumen
`\textgreater{}hue` untuk menunjukkan nilai fungsinya. Untuk fungsi implisit, `xv`
harus berupa fungsi atau ekspresi dari parameter \textdollar{}x\textdollar{} dan \textdollar{}y\textdollar{}, atau
secara alternatif, `xv` dapat berupa matriks nilai.

Euler dapat menandai garis kontur

\end{eulercomment}
\begin{eulerformula}
\[
f(x,y) = c
\]
\end{eulerformula}
\begin{eulercomment}
dari fungsi apa pun.

Untuk menggambar himpunan \textdollar{}f(x,y)=c\textdollar{} bagi satu atau beberapa konstanta
\textdollar{}c\textdollar{}, kamu dapat menggunakan `plot2d()` dengan grafik implisit pada
bidang datar. Parameter untuk \textdollar{}c\textdollar{} adalah `level=c`, di mana \textdollar{}c\textdollar{} dapat
berupa vektor dari garis-garis kontur. Selain itu, skema warna dapat
digambarkan pada latar belakang untuk menunjukkan nilai fungsi pada
setiap titik dalam grafik. Parameter `n` menentukan tingkat kehalusan
(*fineness*) dari grafik tersebut.
\end{eulercomment}
\begin{eulerprompt}
>aspect(1.5); 
>plot2d("x^2+y^2-x*y-x",r=1.5,level=0,contourcolor=red):
\end{eulerprompt}
\eulerimg{17}{Images Use/TugasAplikom4Adi25030630001-136.png}
\begin{eulerprompt}
>expr := "2*x^2+x*y+3*y^4+y"; // define an expression f(x,y)
>plot2d(expr,level=0): // Solutions of f(x,y)=0
\end{eulerprompt}
\eulerimg{17}{Images Use/TugasAplikom4Adi25030630001-137.png}
\begin{eulerprompt}
>plot2d(expr,level=0:0.5:20,>hue,contourcolor=white,n=200): // nice
\end{eulerprompt}
\eulerimg{17}{Images Use/TugasAplikom4Adi25030630001-138.png}
\begin{eulerprompt}
>plot2d(expr,level=0:0.5:20,>hue,>spectral,n=200,grid=4): // nicer
\end{eulerprompt}
\eulerimg{17}{Images Use/TugasAplikom4Adi25030630001-139.png}
\begin{eulercomment}
Opsi ini juga berlaku untuk grafik data (*data plots*). Namun, kamu
harus menentukan rentang (*ranges*) untuk label sumbunya secara
manual.
\end{eulercomment}
\begin{eulerprompt}
>x=-2:0.05:1; y=x'; z=expr(x,y);
>plot2d(z,level=0,a=-1,b=2,c=-2,d=1,>hue):
\end{eulerprompt}
\eulerimg{17}{Images Use/TugasAplikom4Adi25030630001-140.png}
\begin{eulerprompt}
>plot2d("x^3-y^2",>contour,>hue,>spectral):
\end{eulerprompt}
\eulerimg{17}{Images Use/TugasAplikom4Adi25030630001-141.png}
\begin{eulerprompt}
>plot2d("x^3-y^2",level=0,contourwidth=3,>add,contourcolor=red):
\end{eulerprompt}
\eulerimg{17}{Images Use/TugasAplikom4Adi25030630001-142.png}
\begin{eulerprompt}
>z=z+normal(size(z))*0.2;
>plot2d(z,level=0.5,a=-1,b=2,c=-2,d=1):
\end{eulerprompt}
\eulerimg{17}{Images Use/TugasAplikom4Adi25030630001-143.png}
\begin{eulerprompt}
>plot2d(expr,level=[0:0.2:5;0.05:0.2:5.05],color=lightgray):
\end{eulerprompt}
\eulerimg{17}{Images Use/TugasAplikom4Adi25030630001-144.png}
\begin{eulerprompt}
>plot2d("x^2+y^3+x*y",level=1,r=4,n=100):
\end{eulerprompt}
\eulerimg{17}{Images Use/TugasAplikom4Adi25030630001-145.png}
\begin{eulerprompt}
>plot2d("x^2+2*y^2-x*y",level=0:0.1:10,n=100,contourcolor=white,>hue):
\end{eulerprompt}
\eulerimg{17}{Images Use/TugasAplikom4Adi25030630001-146.png}
\begin{eulercomment}
Memungkinkan juga untuk mengisi daerah himpunan

\end{eulercomment}
\begin{eulerformula}
\[
a \le f(x,y) \le b
\]
\end{eulerformula}
\begin{eulercomment}
dengan rentang tingkat (*level range*).

Kamu dapat mengisi daerah nilai untuk fungsi tertentu. Untuk melakukan
hal ini, parameter `level` harus berupa matriks berukuran \textdollar{}2 \textbackslash{}times
n\textdollar{}. Baris pertama berisi batas-batas bawah (*lower bounds*) dan baris
kedua berisi batas-batas atas (*upper bounds*).
\end{eulercomment}
\begin{eulerprompt}
>plot2d(expr,level=[0;1],style="-",color=blue): // 0 <= f(x,y) <= 1
\end{eulerprompt}
\eulerimg{17}{Images Use/TugasAplikom4Adi25030630001-148.png}
\begin{eulercomment}
Grafik implisit (*implicit plots*) juga dapat menampilkan rentang
tingkat (*ranges of levels*). Untuk melakukannya, `level` harus berupa
matriks berukuran \textdollar{}2 \textbackslash{}times n\textdollar{} dari interval tingkat, di mana baris
pertama berisi nilai awal dan baris kedua berisi nilai akhir dari
setiap interval. Secara alternatif, sebuah vektor baris sederhana
dapat digunakan untuk `level`, dan parameter `dl` digunakan untuk
memperluas nilai-nilai `level` tersebut menjadi interval.
\end{eulercomment}
\begin{eulerprompt}
>plot2d("x^4+y^4",r=1.5,level=[0;1],color=blue,style="/"):
\end{eulerprompt}
\eulerimg{17}{Images Use/TugasAplikom4Adi25030630001-149.png}
\begin{eulerprompt}
>plot2d("x^2+y^3+x*y",level=[0,2,4;1,3,5],style="/",r=2,n=100):
\end{eulerprompt}
\eulerimg{17}{Images Use/TugasAplikom4Adi25030630001-150.png}
\begin{eulerprompt}
>plot2d("x^2+y^3+x*y",level=-10:20,r=2,style="-",dl=0.1,n=100):
\end{eulerprompt}
\eulerimg{17}{Images Use/TugasAplikom4Adi25030630001-151.png}
\begin{eulerprompt}
>plot2d("sin(x)*cos(y)",r=pi,>hue,>levels,n=100):
\end{eulerprompt}
\eulerimg{17}{Images Use/TugasAplikom4Adi25030630001-152.png}
\begin{eulercomment}
Dimungkinkan juga untuk menandai suatu daerah

\end{eulercomment}
\begin{eulerformula}
\[
a \le f(x,y) \le b.
\]
\end{eulerformula}
\begin{eulercomment}
Hal ini dilakukan dengan menambahkan parameter `level` yang terdiri
dari dua baris.
\end{eulercomment}
\begin{eulerprompt}
>plot2d("(x^2+y^2-1)^3-x^2*y^3",r=1.3, ...
>  style="#",color=red,<outline, ...
>  level=[-2;0],n=100):
\end{eulerprompt}
\eulerimg{17}{Images Use/TugasAplikom4Adi25030630001-154.png}
\begin{eulercomment}
Memungkinkan juga untuk menentukan nilai tingkat (level) tertentu.
Sebagai contoh, kita dapat menggambar grafik dari solusi suatu
persamaan seperti

\end{eulercomment}
\begin{eulerformula}
\[
x^3-xy+x^2y^2=6
\]
\end{eulerformula}
\begin{eulerprompt}
>plot2d("x^3-x*y+x^2*y^2",r=6,level=1,n=100):
\end{eulerprompt}
\eulerimg{17}{Images Use/TugasAplikom4Adi25030630001-156.png}
\begin{eulerprompt}
>function starplot1 (v, style="/", color=green, lab=none) ...
\end{eulerprompt}
\begin{eulerudf}
    if !holding() then clg; endif;
    w=window(); window(0,0,1024,1024);
    h=holding(1);
    r=max(abs(v))*1.2;
    setplot(-r,r,-r,r);
    n=cols(v); t=linspace(0,2pi,n);
    v=v|v[1]; c=v*cos(t); s=v*sin(t);
    cl=barcolor(color); st=barstyle(style);
    loop 1 to n
      polygon([0,c[#],c[#+1]],[0,s[#],s[#+1]],1);
      if lab!=none then
        rlab=v[#]+r*0.1;
        \{col,row\}=toscreen(cos(t[#])*rlab,sin(t[#])*rlab);
        ctext(""+lab[#],col,row-textheight()/2);
      endif;
    end;
    barcolor(cl); barstyle(st);
    holding(h);
    window(w);
  endfunction
\end{eulerudf}
\begin{eulercomment}
Tidak ada kisi-kisi (*grid*) maupun tanda centang sumbu (*axis ticks*)
di sini. Selain itu, kita menggunakan seluruh jendela (*full window*)
untuk plot tersebut.

Kita memanggil fungsi `reset()` sebelum menguji plot ini untuk
mengembalikan pengaturan bawaan grafis (*graphics defaults*). Hal ini
tidak wajib dilakukan jika kamu sudah yakin bahwa plot tersebut
bekerja dengan baik.
\end{eulercomment}
\begin{eulerprompt}
>reset; starplot1(normal(1,10)+5,color=red,lab=1:10):
\end{eulerprompt}
\eulerimg{27}{Images Use/TugasAplikom4Adi25030630001-157.png}
\begin{eulercomment}
Terkadang, kamu mungkin ingin menggambar sesuatu yang hampir bisa
dilakukan oleh `plot2d`, namun tidak didukung secara langsung.

Pada fungsi berikut, kita membuat plot impuls logaritmik (*logarithmic
impulse plot*). Fungsi `plot2d` memang mampu membuat plot logaritmik,
namun tidak untuk batang impuls (*impulse bars*).
\end{eulercomment}
\begin{eulerprompt}
>function logimpulseplot1 (x,y) ...
\end{eulerprompt}
\begin{eulerudf}
    \{x0,y0\}=makeimpulse(x,log(y)/log(10));
    plot2d(x0,y0,>bar,grid=0);
    h=holding(1);
    frame();
    xgrid(ticks(x));
    p=plot();
    for i=-10 to 10;
      if i<=p[4] and i>=p[3] then
         ygrid(i,yt="10^"+i);
      endif;
    end;
    holding(h);
  endfunction
\end{eulerudf}
\begin{eulercomment}
Mari kita uji dengan nilai-nilai yang berdistribusi eksponensial.
\end{eulercomment}
\begin{eulerprompt}
>aspect(1.5); x=1:10; y=-log(random(size(x)))*200; ...
>logimpulseplot1(x,y):
\end{eulerprompt}
\eulerimg{17}{Images Use/TugasAplikom4Adi25030630001-158.png}
\begin{eulercomment}
Mari kita buat animasi kurva 2D menggunakan grafik langsung (*direct
plots*). Perintah `plot(x,y)` secara sederhana menggambar sebuah kurva
ke dalam jendela grafik. Fungsi `setplot(a,b,c,d)` digunakan untuk
mengatur batas-batas jendela tersebut.

Fungsi `wait(0)` memaksa grafik untuk langsung ditampilkan pada
jendela grafis. Jika tidak, proses menggambar ulang (*redraw*) hanya
akan terjadi dalam interval waktu yang sporadis/jarang.
\end{eulercomment}
\begin{eulerprompt}
>function animliss (n,m) ...
\end{eulerprompt}
\begin{eulerudf}
  t=linspace(0,2pi,500);
  f=0;
  c=framecolor(0);
  l=linewidth(2);
  setplot(-1,1,-1,1);
  repeat
    clg;
    plot(sin(n*t),cos(m*t+f));
    wait(0);
    if testkey() then break; endif;
    f=f+0.02;
  end;
  framecolor(c);
  linewidth(l);
  endfunction
\end{eulerudf}
\begin{eulercomment}
Tekan sembarang tombol untuk menghentikan animasi ini.
\end{eulercomment}
\begin{eulerprompt}
>animliss(2,3); // lihat hasilnya, jika sudah puas, tekan ENTER
\end{eulerprompt}
\eulerheading{Logarithmic Plots}
\begin{eulercomment}
EMT menggunakan parameter logplot untuk skala logaritmik.\\
Grafik logaritmik dapat digambar menggunakan skala logaritmik pada
sumbu \textdollar{}y\textdollar{} dengan logplot=1, skala logaritmik pada sumbu \textdollar{}x\textdollar{} dan \textdollar{}y\textdollar{}
sekaligus dengan logplot=2, atau skala logaritmik hanya pada sumbu \textdollar{}x\textdollar{}
dengan logplot=3.

\end{eulercomment}
\begin{eulerttcomment}
 - logplot=1: sumbu $y$ logaritmik
 - logplot=2: sumbu $x$ dan $y$ 
 - logaritmiklogplot=3: sumbu $x$ logaritmik
\end{eulerttcomment}
\begin{eulerprompt}
>plot2d("exp(x^3-x)*x^2",1,5,logplot=1):
\end{eulerprompt}
\eulerimg{17}{Images Use/TugasAplikom4Adi25030630001-159.png}
\begin{eulerprompt}
>plot2d("exp(x+sin(x))",0,100,logplot=1):
\end{eulerprompt}
\eulerimg{17}{Images Use/TugasAplikom4Adi25030630001-160.png}
\begin{eulerprompt}
>plot2d("exp(x+sin(x))",10,100,logplot=2):
\end{eulerprompt}
\eulerimg{17}{Images Use/TugasAplikom4Adi25030630001-161.png}
\begin{eulerprompt}
>plot2d("gamma(x)",1,10,logplot=1):
\end{eulerprompt}
\eulerimg{17}{Images Use/TugasAplikom4Adi25030630001-162.png}
\begin{eulerprompt}
>plot2d("log(x*(2+sin(x/100)))",10,1000,logplot=3):
\end{eulerprompt}
\eulerimg{17}{Images Use/TugasAplikom4Adi25030630001-163.png}
\begin{eulercomment}
Hal ini juga berlaku untuk grafik data (*data plots*).
\end{eulercomment}
\begin{eulerprompt}
>x=10^(1:20); y=x^2-x;
>plot2d(x,y,logplot=2):
\end{eulerprompt}
\eulerimg{17}{Images Use/TugasAplikom4Adi25030630001-164.png}
\begin{eulerprompt}
>                                                      
\end{eulerprompt}
\eulerheading{Contoh Penyelesaian Masalah Menggambar Kurva/Plot 2D}
\begin{eulercomment}
\end{eulercomment}
\eulersubheading{1. Materi Aljabar (Titik Potong Kurva Parabola dan Garis Linear)}
\begin{eulercomment}
Soal 1: Gambarkan grafik fungsi kuadrat dan fungsi linear berikut pada satu
bidang koordinat yang sama, serta tentukan dan tandai titik potong kedua kurva tersebut!

\end{eulercomment}
\begin{eulerformula}
\[
f(x) = x^2 - 2x - 3
\]
\end{eulerformula}
\begin{eulerformula}
\[
g(x) = x + 1
\]
\end{eulerformula}
\begin{eulercomment}
\end{eulercomment}
\begin{eulerprompt}
>aspect(1.5);
>plot2d(["x^2 - 2*x - 3", "x + 1"], a=-3, b=5, c=-5, d=8, ...
>  color=[blue, red], style=["-", "--"], grid=6, ...
>  xl="x", yl="y", title="Titik Potong Parabola dan Garis"); ...
>plot2d([-1, 4], [0, 5], >points, style="o#", color=black, >add); ...
>label("(-1, 0)", -1, 0, pos="ul", offset=10); ...
>label("(4, 5)", 4, 5, pos="ul", offset=10); ...
>labelbox(["f(x) = x^2 - 2x - 3", "g(x) = x + 1"], styles=["-", "--"], colors=[blue, red], x=0.1, y=0.9, >left):
\end{eulerprompt}
\eulerimg{17}{Images Use/TugasAplikom4Adi25030630001-165.png}
\begin{eulercomment}
Penjelasan hasil output:
Grafik menampilkan kurva parabola terbuka ke atas (biru) yang berpotongan dengan
garis lurus bergradien positif (merah putus-putus) di dua titik, yaitu (-1, 0) dan (4, 5).
Titik potong ini sesuai dengan penyelesaian aljabar x\textasciicircum{}2 - 2x - 3 = x + 1 yang
menghasilkan akar-akar x = -1 dan x = 4.
\end{eulercomment}
\eulersubheading{2. Materi Kalkulus (Kurva Fungsi dan Turunan Pertamanya)}
\begin{eulercomment}
Soal 2: Diketahui fungsi polinomial berderajat tiga berikut:

\end{eulercomment}
\begin{eulerformula}
\[
f(x) = x^3 - 3x + 1
\]
\end{eulerformula}
\begin{eulercomment}
Gambarkan grafik f(x) beserta turunan pertamanya f'(x) pada selang [-2.5, 2.5],
serta tandai titik ekstrem lokal (maksimum dan minimum lokal) pada kurva f(x)!

\end{eulercomment}
\begin{eulerprompt}
>function f(x) &= x^3 - 3*x + 1;
>plot2d(&f(x), a=-2.5, b=2.5, c=-4, d=5, color=black, thickness=2, grid=6, ...
>  xl="x", yl="y", title="Fungsi Kubik dan Turunannya"); ...
>plot2d(&diff(f(x), x), >add, color=blue, style="--"); ...
>plot2d([-1, 1], [f(-1), f(1)], >points, style="o#", color=red, >add); ...
>label("Maks Lokal (-1, 3)", -1, f(-1), pos="bc", offset=10); ...
>label("Min Lokal (1, -1)", 1, f(1), pos="tc", offset=10); ...
>labelbox(["f(x)", "f'(x)"], styles=["-", "--"], colors=[black, blue]):
\end{eulerprompt}
\eulerimg{17}{Images Use/TugasAplikom4Adi25030630001-166.png}
\begin{eulercomment}
Penjelasan hasil output:
Kurva hitam adalah grafik fungsi kubik f(x) = x\textasciicircum{}3 - 3x + 1, sedangkan kurva biru
putus-putus adalah turunan pertamanya f'(x) = 3x\textasciicircum{}2 - 3. Terlihat jelas bahwa
saat grafik turunan f'(x) memotong sumbu-x (bernilai nol) di x = -1 dan x = 1,
kurva f(x) mencapai titik stasioner, yaitu titik maksimum lokal di (-1, 3)
dan titik minimum lokal di (1, -1).
\end{eulercomment}
\eulersubheading{3. Materi Trigonometri (Perbandingan Amplitudo dan Periode Fungsi Sinus)}
\begin{eulercomment}
Soal 3: Gambarkan grafik kedua fungsi trigonometri berikut pada interval

\end{eulercomment}
\begin{eulerformula}
\[
0 \le x \le 2\pi
\]
\end{eulerformula}
\begin{eulercomment}
untuk mengamati perbedaan amplitudo dan frekuensinya:

\end{eulercomment}
\begin{eulerformula}
\[
y_1 = \sin(x) \quad \text{dan} \quad y_2 = 2\sin(3x)
\]
\end{eulerformula}
\begin{eulercomment}
\end{eulercomment}
\begin{eulerprompt}
>plot2d(["sin(x)", "2*sin(3*x)"], a=0, b=2*pi, c=-2.5, d=2.5, ...
>  color=[blue, red], style=["-", "-."], grid=6, ...
>  xl="x (radian)", yl="y", title="Perbandingan Amplitudo dan Periode Sinus"); ...
>labelbox(["sin(x)", "2 sin(3x)"], styles=["-", "-."], colors=[blue, red]):
\end{eulerprompt}
\eulerimg{17}{Images Use/TugasAplikom4Adi25030630001-167.png}
\begin{eulercomment}
Penjelasan hasil output:
Grafik biru menunjukkan fungsi dasar sin(x) dengan amplitudo 1 dan 1 gelombang penuh
(periode 2*pi) pada selang [0, 2*pi]. Grafik merah menunjukkan fungsi 2*sin(3x) yang
memiliki amplitudo dua kali lebih tinggi (bernilai antara -2 sampai 2) dan frekuensi
tiga kali lebih rapat sehingga terbentuk 3 gelombang penuh (periode 2*pi/3).
\end{eulercomment}
\eulersubheading{4. Materi Geometri Analitik / Fungsi Implisit (Irisan Kerucut: Elips dan Hiperbola)}
\begin{eulercomment}
Soal 4: Gambarkan kurva implisit dari persamaan elips dan hiperbola berikut pada
bidang koordinat yang sama dengan rasio aspek persegi (square):

\end{eulercomment}
\begin{eulerformula}
\[
\frac{x^2}{9} + \frac{y^2}{4} = 1 \quad \text{dan} \quad x^2 - y^2 = 1
\]
\end{eulerformula}
\begin{eulercomment}
\end{eulercomment}
\begin{eulerprompt}
>reset;
>plot2d("x^2/9 + y^2/4 - 1", r=4, level=0, contourcolor=blue, contourwidth=2, >square, grid=6, ...
>  xl="x", yl="y", title="Kurva Implisit Elips dan Hiperbola"); ...
>plot2d("x^2 - y^2 - 1", level=0, contourcolor=red, contourwidth=2, >add):
\end{eulerprompt}
\eulerimg{27}{Images Use/TugasAplikom4Adi25030630001-168.png}
\begin{eulercomment}
Penjelasan hasil output:
Dengan menggunakan parameter level=0 pada fungsi implisit, EMT menggambar dua kurva
irisan kerucut secara akurat. Kurva biru membentuk elips horizontal yang berpusat di (0,0)
dengan sumbu semi-mayor 3 (memotong sumbu-x di -3 dan 3) dan sumbu semi-minor 2
(memotong sumbu-y di -2 dan 2). Kurva merah membentuk hiperbola terbuka ke kiri dan kanan
dengan titik puncak di (-1, 0) dan (1, 0) yang memotong elips di empat titik simetris.
\end{eulercomment}
\begin{eulerprompt}
>  
\end{eulerprompt}
\eulerheading{Rujukan Lengkap Fungsi plot2d()}
\begin{eulercomment}
\end{eulercomment}
\begin{eulerttcomment}
  function plot2d (xv, yv, btest, a, b, c, d, xmin, xmax, r, n,  ..
  logplot, grid, frame, framecolor, square, color, thickness, style, ..
  auto, add, user, delta, points, addpoints, pointstyle, bar, histogram,  ..
  distribution, even, steps, own, adaptive, hue, level, contour,  ..
  nc, filled, fillcolor, outline, title, xl, yl, maps, contourcolor, ..
  contourwidth, ticks, margin, clipping, cx, cy, insimg, spectral,  ..
  cgrid, vertical, smaller, dl, niveau, levels)
\end{eulerttcomment}
\begin{eulercomment}
Multipurpose plot function for plots in the plane (2D plots). This function can do
plots of functions of one variables, data plots, curves in the plane, bar plots, grids
of complex numbers, and implicit plots of functions of two variables.

Parameters
\\
x,y       : equations, functions or data vectors\\
a,b,c,d   : Plot area (default a=-2,b=2)\\
r         : if r is set, then a=cx-r, b=cx+r, c=cy-r, d=cy+r\\
\end{eulercomment}
\begin{eulerttcomment}
            r can be a vector [rx,ry] or a vector [rx1,rx2,ry1,ry2].
\end{eulerttcomment}
\begin{eulercomment}
xmin,xmax : range of the parameter for curves\\
auto      : Determine y-range automatically (default)\\
square    : if true, try to keep square x-y-ranges\\
n         : number of intervals (default is adaptive)\\
grid      : 0 = no grid and labels,\\
\end{eulercomment}
\begin{eulerttcomment}
            1 = axis only,
            2 = normal grid (see below for the number of grid lines)
            3 = inside axis
            4 = no grid
            5 = full grid including margin
            6 = ticks at the frame
            7 = axis only
            8 = axis only, sub-ticks
\end{eulerttcomment}
\begin{eulercomment}
frame     : 0 = no frame\\
framecolor: color of the frame and the grid\\
margin    : number between 0 and 0.4 for the margin around the plot\\
color     : Color of curves. If this is a vector of colors,\\
\end{eulercomment}
\begin{eulerttcomment}
            it will be used for each row of a matrix of plots. In the case of
            point plots, it should be a column vector. If a row vector or a
            full matrix of colors is used for point plots, it will be used for
            each data point.
\end{eulerttcomment}
\begin{eulercomment}
thickness : line thickness for curves\\
\end{eulercomment}
\begin{eulerttcomment}
            This value can be smaller than 1 for very thin lines.
\end{eulerttcomment}
\begin{eulercomment}
style     : Plot style for lines, markers, and fills.\\
\end{eulercomment}
\begin{eulerttcomment}
            For points use
            "[]", "<>", ".", "..", "...",
            "*", "+", "|", "-", "o"
            "[]#", "<>#", "o#" (filled shapes)
            "[]w", "<>w", "ow" (non-transparent)
            For lines use
            "-", "--", "-.", ".", ".-.", "-.-", "->"
            For filled polygons or bar plots use
            "#", "#O", "O", "/", "\(\backslash\)", "\(\backslash\)/",
            "+", "|", "-", "t"
\end{eulerttcomment}
\begin{eulercomment}
points    : plot single points instead of line segments\\
addpoints : if true, plots line segments and points\\
add       : add the plot to the existing plot\\
user      : enable user interaction for functions\\
delta     : step size for user interaction\\
bar       : bar plot (x are the interval bounds, y the interval values)\\
histogram : plots the frequencies of x in n subintervals\\
distribution=n : plots the distribution of x with n subintervals\\
even      : use inter values for automatic histograms.\\
steps     : plots the function as a step function (steps=1,2)\\
adaptive  : use adaptive plots (n is the minimal number of steps)\\
level     : plot level lines of an implicit function of two variables\\
outline   : draws boundary of level ranges.
\\
If the level value is a 2xn matrix, ranges of levels will be drawn\\
in the color using the given fill style. If outline is true, it\\
will be drawn in the contour color. Using this feature, regions of\\
f(x,y) between limits can be marked.
\\
hue       : add hue color to the level plot to indicate the function\\
\end{eulercomment}
\begin{eulerttcomment}
            value
\end{eulerttcomment}
\begin{eulercomment}
contour   : Use level plot with automatic levels\\
nc        : number of automatic level lines\\
title     : plot title (default "")\\
xl, yl    : labels for the x- and y-axis\\
smaller   : if \textgreater{}0, there will be more space to the left for labels.\\
vertical  :\\
\end{eulercomment}
\begin{eulerttcomment}
  Turns vertical labels on or off. This changes the global variable
  verticallabels locally for one plot. The value 1 sets only vertical
  text, the value 2 uses vertical numerical labels on the y axis.
\end{eulerttcomment}
\begin{eulercomment}
filled    : fill the plot of a curve\\
fillcolor : fill color for bar and filled curves\\
outline   : boundary for filled polygons\\
logplot   : set logarithmic plots\\
\end{eulercomment}
\begin{eulerttcomment}
            1 = logplot in y,
            2 = logplot in xy,
            3 = logplot in x
\end{eulerttcomment}
\begin{eulercomment}
own       :\\
\end{eulercomment}
\begin{eulerttcomment}
  A string, which points to an own plot routine. With >user, you get
  the same user interaction as in plot2d. The range will be set
  before each call to your function.
\end{eulerttcomment}
\begin{eulercomment}
maps      : map expressions (0 is faster), functions are always mapped.\\
contourcolor : color of contour lines\\
contourwidth : width of contour lines\\
clipping  : toggles the clipping (default is true)\\
title     :\\
\end{eulercomment}
\begin{eulerttcomment}
  This can be used to describe the plot. The title will appear above
  the plot. Moreover, a label for the x and y axis can be added with
  xl="string" or yl="string". Other labels can be added with the
  functions label() or labelbox(). The title can be a unicode
  string or an image of a Latex formula.
\end{eulerttcomment}
\begin{eulercomment}
cgrid     :\\
\end{eulercomment}
\begin{eulerttcomment}
  Determines the number of grid lines for plots of complex grids.
  Should be a divisor of the the matrix size minus 1 (number of
  subintervals). cgrid can be a vector [cx,cy].
\end{eulerttcomment}
\begin{eulercomment}

Overview

The function can plot

- expressions, call collections or functions of one variable,\\
- parametric curves,\\
- x data against y data,\\
- implicit functions,\\
- bar plots,\\
- complex grids,\\
- polygons.

If a function or expression for xv is given, plot2d() will compute\\
values in the given range using the function or expression. The\\
expression must be an expression in the variable x. The range must\\
be defined in the parameters a and b unless the default range\\
[-2,2] should be used. The y-range will be computed automatically,\\
unless c and d are specified, or a radius r, which yields the range\\
[-r,r] for x and y. For plots of functions, plot2d will use an\\
adaptive evaluation of the function by default. To speed up the\\
plot for complicated functions, switch this off with \textless{}adaptive, and\\
optionally decrease the number of intervals n. Moreover, plot2d()\\
will by default use mapping. I.e., it will compute the plot element\\
for element. If your expression or your functions can handle a\\
vector x, you can switch that off with \textless{}maps for faster evaluation.

Note that adaptive plots are always computed element for element. \\
If functions or expressions for both xv and for yv are specified,\\
plot2d() will compute a curve with the xv values as x-coordinates\\
and the yv values as y-coordinates. In this case, a range should be\\
defined for the parameter using xmin, xmax. Expressions contained\\
in strings must always be expressions in the parameter variable x.
\end{eulercomment}
\end{eulernotebook}
\end{document}


% ====================================================================



% ====================================================================
\end{document}